\chapter{Synthesis, Reference Architectures \& Future Directions}
\label{ch:synthesis}

\section*{Executive Summary}

Every chapter so far has taken one API concern --- the wire protocol, a
framework, an error model, a scaling tactic, a governance discipline --- and
treated it with the depth it deserves. Real systems, however, are never about
one concern. A production API platform is a \emph{composition}: HTTP semantics
(Chapter~\ref{ch:http_wire}), REST resource design
(Chapter~\ref{ch:rest_style}), GraphQL selection sets
(Chapter~\ref{ch:graphql}), gRPC streaming (Chapter~\ref{ch:grpc}), event
delivery (Chapter~\ref{ch:async_apis}), and the operational disciplines of
security, observability, testing, and lifecycle governance
(Chapters~\ref{ch:api_security}--\ref{ch:lifecycle}) all coexist inside one
coherent architecture. This final chapter is about that composition.

The chapter delivers four capstone artifacts. First, the \textbf{Master
Decision Card}: a rigorous, constraint-driven selection framework for choosing
between REST, GraphQL, gRPC, event-driven/webhook, and WebSocket/SSE
interaction profiles --- expressed as decision tables, a decision tree, and
worked scenarios, and then \emph{encoded as executable software}. Second, the
\textbf{Northwind Robotics API Platform reference architecture}: one
integrated design that uses every concept in the book, with per-layer
commentary mapping each architectural element back to the chapter that taught
it. Third, the \textbf{Grand Capstone}: a complete, buildable specification of
that platform --- scope, domains, API surfaces, non-functional requirements, a
four-milestone delivery plan, acceptance criteria, and an evaluation rubric.
Fourth, an honest assessment of the \textbf{emerging frontiers} --- HTTP/3 and
QUIC, AI-agent-consumed APIs and the Model Context Protocol, AI-assisted
design, edge-native deployment, post-quantum TLS, and security automation ---
so that you can separate deployable reality from conference-stage enthusiasm.

The chapter closes the loop on your own development: a cross-cutting
synthesis of the book's failure taxonomy, one-page-per-module summary tables,
career pathways from junior to staff API engineer mapped to chapters, a
consolidated bibliography, and five integrative capstone projects. The code in
this chapter is deliberately integrative: a protocol-selection advisor with a
full rationale trace, a single FastAPI application that exposes REST, GraphQL,
and SSE surfaces with RFC~9457 errors, idempotency keys, and rate-limit
headers, and an architecture conformance checker that turns the book's
governance rules into executable policy.

\section{Learning Objectives \& Prerequisites}

After completing this chapter, you will be able to:

\begin{itemize}
  \item Select an API interaction profile (REST, GraphQL, gRPC,
        event-driven, or streaming) from stated constraints, and defend the
        selection with an explicit, auditable rationale.
  \item Assemble a complete reference architecture --- edge, gateway,
        services, data, observability, and governance --- and justify every
        layer against the non-functional requirements it serves.
  \item Specify a multi-protocol API platform to the standard a staff
        engineer would expect: SLOs, security posture, rate tiers, delivery
        milestones, and measurable acceptance criteria.
  \item Diagnose which chapter's discipline applies to any concrete failure
        mode observed in production.
  \item Assess emerging API technologies (HTTP/3, MCP, edge-native compute,
        post-quantum TLS) with a sober adopt/trial/wait classification.
  \item Map your own competency growth from junior to staff engineer onto
        concrete, chapter-anchored skills.
  \item Encode architectural governance as executable policy that runs
        offline in CI, and prove platform surfaces work together with
        in-process smoke tests.
\end{itemize}

\subsection*{Prerequisites}

This is the synthesis chapter: it assumes the whole book. You should be
comfortable with HTTP message semantics and connection management
(Chapter~\ref{ch:http_wire}), client-side consumption patterns
(Chapter~\ref{ch:consuming_python}), REST maturity and resource modeling
(Chapter~\ref{ch:rest_style}), serialization contracts
(Chapter~\ref{ch:serialization}), building services with FastAPI
(Chapter~\ref{ch:fastapi}), RFC~9457 error design
(Chapter~\ref{ch:problem_details}), pagination and versioning
(Chapters~\ref{ch:pagination}--\ref{ch:versioning}), idempotency and caching
(Chapter~\ref{ch:idempotency}), asynchronous and event-driven APIs
(Chapter~\ref{ch:async_apis}), GraphQL and gRPC
(Chapters~\ref{ch:graphql}--\ref{ch:grpc}), rate limiting
(Chapter~\ref{ch:rate_limiting}), microservice resilience
(Chapter~\ref{ch:microservices}), security engineering
(Chapter~\ref{ch:api_security}), deployment (Chapter~\ref{ch:deployment}),
observability (Chapter~\ref{ch:observability}), testing
(Chapter~\ref{ch:testing}), and lifecycle governance
(Chapter~\ref{ch:lifecycle}). Sections below reference these chapters
liberally rather than re-teaching them.

\section{Deep Technical Mechanics \& Architectural Concepts}

\subsection{The Master Decision Card: Choosing an Interaction Profile}
\label{sec:decision-card}

The most consequential API decision a team makes is rarely \emph{how} to
implement a protocol --- it is \emph{which} protocol to implement. Choosing
poorly is expensive in a way that refactoring cannot fix: the protocol is the
contract your consumers build against, and contracts accrete
(Chapter~\ref{ch:lifecycle}). The industry conversation about this choice is
unfortunately dominated by fashion. The antidote is to treat protocol
selection as \textbf{constraint satisfaction}, not taste.

\begin{definition}[Interaction Profile]
An \emph{interaction profile} is the combination of transport semantics,
message format, delivery model, and contract style through which a consumer
and provider exchange data. This book covers five: \textbf{REST} (resource
oriented, HTTP semantics, Chapter~\ref{ch:rest_style}), \textbf{GraphQL}
(selection sets over a typed graph, Chapter~\ref{ch:graphql}), \textbf{gRPC}
(contract-first RPC over HTTP/2 with Protobuf, Chapter~\ref{ch:grpc}),
\textbf{event-driven} (brokered events or signed webhooks,
Chapter~\ref{ch:async_apis}), and \textbf{streaming} (Server-Sent Events or
WebSockets, Chapter~\ref{ch:async_apis}).
\end{definition}

\subsubsection{The Six Criteria}

The decision card evaluates a candidate interaction against six criteria.
Each criterion is a question with a small set of allowed answers --- which is
precisely what makes the framework encodable as the questionnaire of
Listing~1 (Section~\ref{sec:code-advisor}).

\begin{table}[htbp]
  \centering
  \small
  \begin{tabularx}{\textwidth}{l l X}
    \toprule
    \textbf{Criterion} & \textbf{Question} & \textbf{Why it dominates} \\
    \midrule
    Consumer type & Who calls this API? &
      Browser-based third parties cannot speak raw gRPC; machines do not care
      about human-readable payloads; mobile clients feel every over-fetched
      byte on a metered radio. \\
    Payload shape & Fixed resources or a nested graph? &
      Deep object graphs make REST chatty (the N+1 endpoint problem) and make
      GraphQL's selection sets pay for themselves. \\
    Latency budget & What p99 must you hold? &
      Sub-50\,ms service-to-service budgets favor gRPC's binary framing and
      HTTP/2 multiplexing \cite{rfc9113,grpcdocs}; relaxed budgets tolerate
      store-and-forward asynchrony. \\
    Delivery model & Sync, push, or fire-and-forget? &
      Structural constraint: unary request--response protocols cannot push.
      If the server must initiate, you need SSE/WebSockets or events. \\
    HTTP caching & Must shared caches help? &
      Only REST GETs are addressable, cacheable resources under HTTP
      semantics \cite{rfc9110}; POST-based GraphQL bypasses the CDN layer
      entirely. \\
    Team maturity & Exploratory or disciplined? &
      GraphQL and gRPC pay off only with schema governance discipline; REST
      has the shallowest tooling curve \cite{gehl2021apidesign}. \\
    \bottomrule
  \end{tabularx}
  \caption{The six criteria of the Master Decision Card.}
  \label{tab:decision-criteria}
\end{table}

\subsubsection{The Comparison Table}

Table~\ref{tab:protocol-matrix} condenses twenty chapters of trade-offs into
one grid. Read it as \emph{tendencies}, not laws --- each row has an entire
chapter of nuance behind it.

\begin{table}[htbp]
  \centering
  \small
  \begin{tabularx}{\textwidth}{l X X X X X}
    \toprule
    & \textbf{REST} & \textbf{GraphQL} & \textbf{gRPC} & \textbf{Events/webhooks} & \textbf{SSE/WS} \\
    \midrule
    Primary consumer &
      public, partner &
      mobile, SPA, data &
      internal services &
      machines, partners &
      dashboards, tickers \\
    Delivery &
      sync & sync & sync + streams & async, at-least-once & server push \\
    Payload &
      JSON resources \cite{rfc8259} &
      client-shaped JSON &
      binary Protobuf \cite{protobufdocs} &
      event envelopes &
      text frames / JSON \\
    Caching &
      full HTTP caching \cite{rfc9110} &
      persisted queries only &
      none &
      n/a (push) &
      no-cache streams \\
    Contract &
      OpenAPI \cite{openapi31} &
      SDL schema \cite{graphqlspec} &
      proto files \cite{grpcdocs} &
      AsyncAPI \cite{asyncapi} &
      AsyncAPI \\
    Tooling curve &
      shallow & medium & medium--steep & medium & shallow--medium \\
    Failure mode &
      N+1 hops & N+1 resolvers & coupling via stubs &
      silent event loss & connection storms \\
    Book chapter &
      \ref{ch:rest_style} & \ref{ch:graphql} & \ref{ch:grpc} &
      \ref{ch:async_apis} & \ref{ch:async_apis} \\
    \bottomrule
  \end{tabularx}
  \caption{Interaction profile comparison matrix (tendencies, not laws).}
  \label{tab:protocol-matrix}
\end{table}

\begin{keyconceptbox}
\textbf{Constraints choose protocols; engineers choose implementations.}
If your stated constraints (consumer, payload, latency, delivery, caching,
maturity) do not uniquely determine a profile, your constraints are not yet
stated sharply enough. When two profiles genuinely tie, prefer the one your
organization already operates well: operational maturity beats theoretical
fit.
\end{keyconceptbox}

\subsubsection{The Decision Tree}

Figure~\ref{fig:decision-tree} renders the card as a tree. The delivery model
comes first because it is a \emph{structural} constraint --- no amount of
engineering makes a unary REST endpoint push data to the client.

\begin{figure}[htbp]
\centering
\begin{tikzpicture}[
  font=\small,
  >=stealth,
  node distance=5mm and 10mm,
  qnode/.style={flowstep, fill=blue!5, align=center},
  ansnode/.style={flowstep, fill=green!5, align=center},
  lbl/.style={note, midway}
]
\node[qnode] (d1) {Q4: delivery\\model?};
\node[qnode, below left=14mm and 6mm of d1] (d2) {Q1+Q3: consumer and\\latency budget?};
\node[ansnode, right=18mm of d1] (ev) {Event-driven\\webhooks / broker\\(Ch.~\ref{ch:async_apis})};
\node[ansnode, above right=4mm and 18mm of d1] (st) {SSE / WebSockets\\(Ch.~\ref{ch:async_apis})};
\node[ansnode, below left=6mm and 2mm of d2] (grpc) {gRPC\\(Ch.~\ref{ch:grpc})};
\node[qnode, below right=6mm and 2mm of d2] (d3) {Q2: payload\\shape?};
\node[ansnode, below left=6mm and 2mm of d3] (rest) {REST + OpenAPI\\(Ch.~\ref{ch:rest_style})};
\node[ansnode, below right=6mm and 2mm of d3] (gql) {GraphQL\\(Ch.~\ref{ch:graphql})};

\draw[->] (d1) -- node[lbl, above left] {sync} (d2);
\draw[->] (d1) -- node[lbl, above] {fire-and-forget} (ev);
\draw[->] (d1) -- node[lbl, above] {server push} (st);
\draw[->] (d2) -- node[lbl, above left] {internal + p99 $<$ 50\,ms} (grpc);
\draw[->] (d2) -- node[lbl, above right] {otherwise} (d3);
\draw[->] (d3) -- node[lbl, above left] {fixed resources / caching} (rest);
\draw[->] (d3) -- node[lbl, above right] {nested graph, mobile/SPA} (gql);
\end{tikzpicture}
\caption{The Master Decision Card as a decision tree. Question numbers refer
to Table~\ref{tab:decision-criteria}; delivery is tested first because it is
structural. Edge cases (gRPC-Web for browsers, GraphQL subscriptions over
WebSockets) are refinements, not new branches.}
\label{fig:decision-tree}
\end{figure}

\subsubsection{Worked Scenarios}

The framework earns its keep on realistic, messy scenarios. Each of the five
below is encoded verbatim as a canonical test case in the advisor of
Section~\ref{sec:code-advisor}, so the prose and the executable policy cannot
drift apart.

\begin{examplebox}
\textbf{Scenario A --- Public SaaS API.} A vendor exposes its product catalog
to unknown third-party developers. Consumer: public. Payload: fixed
resources. Latency: relaxed (hundreds of milliseconds fine). Delivery:
synchronous. Caching: catalog reads must be CDN-cacheable. Team: disciplined.
\textbf{Recommendation: REST + OpenAPI.} Third parties expect plain HTTPS and
JSON they can explore with \texttt{curl}; shared caching
\cite{rfc9110} absorbs read amplification; OpenAPI gives the portal its
source of truth \cite{openapi31}. GraphQL adds support burden with no
caching upside; gRPC is unreachable from browsers without a proxy layer.
\end{examplebox}

\textbf{Scenario B --- Internal microservice mesh.} Forty services exchange
inventory and pricing calls inside one trust boundary. Consumer: internal.
Latency: extreme (p99 $<$ 50\,ms budgets on hot paths). Delivery: sync, with
some server-streaming feeds. \textbf{Recommendation: gRPC.} Strongly typed
contracts, code generation, HTTP/2 multiplexing, and binary framing
\cite{rfc9113,grpcdocs} dominate at this latency budget, and the team
controls both ends of every call.

\textbf{Scenario C --- Mobile app backend.} A field-technician app renders
deeply nested work-order screens over metered connections. Consumer: mobile.
Payload: nested graph. Latency: tight. Caching: personalized, so no shared
cache. \textbf{Recommendation: GraphQL.} One round trip per screen with
exactly the fields rendered \cite{graphqlspec}; an aggregating BFF-style REST
layer is the acceptable fallback when GraphQL governance is not yet in place.

\textbf{Scenario D --- IoT telemetry.} Three thousand robots stream sensor
readings. Consumer: machine. Payload: high-volume streams. Delivery:
fire-and-forget. \textbf{Recommendation: event-driven} ingestion to a broker
with an event envelope (\texttt{event\_id}, \texttt{occurred\_at});
gRPC bidirectional streaming is the runner-up when devices speak HTTP/2 and
the operator wants backpressure-aware streams
(Chapter~\ref{ch:async_apis}).

\textbf{Scenario E --- Partner integrations.} Business partners must learn
within seconds that an order shipped. Consumer: public, delivery:
fire-and-forget to \emph{their} endpoints. \textbf{Recommendation: signed
webhooks} with delivery logs and replay, per
Chapter~\ref{ch:async_apis}. Polling a REST endpoint is the graceful
degradation path, not the primary design.

\begin{warningbox}
\textbf{The decision card is a floor, not a ceiling.} It prevents
\emph{indefensible} choices (public gRPC to browsers, unary APIs for push,
uncached GraphQL for a read-heavy public catalog). It does not certify that
the chosen design is good --- that still requires the discipline of the
relevant chapter.
\end{warningbox}

\subsection{The Northwind Robotics Reference Architecture}
\label{sec:ref-arch}

We now integrate the entire book into one design: the \textbf{Northwind
Robotics API Platform}, a synthetic-but-realistic platform for a robotics
manufacturer. Northwind sells robots and spare parts (catalog, orders), runs
a connected fleet (telemetry), and integrates with resellers (partner APIs).
Figure~\ref{fig:ref-arch} shows the full picture; the commentary after it
maps every layer to the chapters that justify it.

\begin{figure}[htbp]
\centering
\begin{tikzpicture}[
  font=\small,
  >=stealth,
  node distance=4mm and 6mm,
  pbox/.style={flowstep, fill=white, align=center},
  layerbox/.style={draw=packtgray, dashed, inner sep=6pt},
  lbl/.style={note}
]
% --- consumers ---
\node[pbox] (spa) {Partner\\portals};
\node[pbox, right=of spa] (mobile) {Mobile\\app};
\node[pbox, right=of mobile] (web) {Web\\storefront};
\node[pbox, right=of web] (fleet) {Robot\\fleet};

% --- edge ---
\node[pbox, below=8mm of mobile, fill=blue!5] (edge) {CDN + WAF\\ TLS 1.3 termination \cite{rfc8446}\\ static cache \cite{rfc9110}};

% --- gateway ---
\node[pbox, below=6mm of edge, fill=blue!5] (gw) {API gateway\\ auth offload \cite{rfc6749,rfc7519}\\ rate tiers (Ch.~\ref{ch:rate_limiting})\\ request validation};

% --- services ---
\node[pbox, below left=8mm and 16mm of gw] (rest) {Catalog + orders\\REST / FastAPI\\(Ch.~\ref{ch:fastapi})};
\node[pbox, below=8mm of gw] (gql) {GraphQL\\read model\\(Ch.~\ref{ch:graphql})};
\node[pbox, below right=8mm and 16mm of gw] (grpc) {Inventory\\gRPC internal\\(Ch.~\ref{ch:grpc})};
\node[pbox, right=8mm of grpc] (push) {Webhooks + SSE\\dispatcher\\(Ch.~\ref{ch:async_apis})};

% --- data ---
\node[pbox, below=10mm of gql] (db) {PostgreSQL\\+ outbox table};
\node[pbox, right=10mm of db] (broker) {Event broker\\(telemetry bus)};
\node[pbox, left=10mm of db] (cache) {Redis\\rate limits, cache};

% --- cross-cutting ---
\node[pbox, below=8mm of db, fill=yellow!5] (otel) {Observability: OpenTelemetry traces, metrics, logs \cite{opentelemetrydocs} (Ch.~\ref{ch:observability})};
\node[pbox, below=6mm of otel, fill=yellow!5] (gov) {Governance: linted OpenAPI/AsyncAPI specs, developer portal, conformance CI \cite{openapi31,asyncapi} (Ch.~\ref{ch:lifecycle})};

% --- edges ---
\draw[->] (spa) -- (edge);
\draw[->] (mobile) -- (edge);
\draw[->] (web) -- (edge);
\draw[->] (fleet) -- (broker);
\draw[->] (edge) -- (gw);
\draw[->] (gw) -- (rest);
\draw[->] (gw) -- (gql);
\draw[->] (gw) -- (grpc);
\draw[->] (rest) -- (db);
\draw[->] (gql) -- (db);
\draw[->] (grpc) -- (db);
\draw[->] (db) -- node[lbl, above] {outbox relay} (broker);
\draw[->] (broker) -- (push);
\draw[->] (push) -- node[lbl, right] {signed events} (spa);
\draw[->] (push) -- node[lbl, right] {SSE} (mobile);
\draw[<->] (gw) -- (cache);
\draw[->, dashed] (rest) -- (otel);
\draw[->, dashed] (gql) -- (otel);
\draw[->, dashed] (grpc) -- (otel);
\draw[->, dashed] (gw) -- (otel);
\draw[->, dashed] (gov) -- (gw);

% --- layer labels ---
\node[lbl, left=2mm of edge] {Edge};
\node[lbl, left=2mm of gw] {Gateway};
\node[lbl, left=2mm of gql, xshift=-6mm] {Services};
\node[lbl, left=2mm of db] {Data};
\end{tikzpicture}
\caption{The Northwind Robotics API Platform reference architecture. Solid
arrows are request/event flow; dashed arrows are control-plane and telemetry
flow. Every box names the chapter that teaches it.}
\label{fig:ref-arch}
\end{figure}

\subsubsection{Layer-by-Layer Commentary}

\textbf{Consumers.} Four consumer classes drive four interaction profiles,
exactly as the decision card predicts: partners (REST + webhooks), mobile
(GraphQL + SSE), the web storefront (REST with CDN caching), and the robot
fleet (event-driven ingestion). One platform, four profiles --- chosen per
consumer, not per fashion.

\textbf{Edge.} A CDN terminates TLS~1.3 \cite{rfc8446}, serves cached catalog
GETs, and absorbs cheap abuse; a WAF filters known-bad signatures before they
cost compute. This is where HTTP semantics pay rent: correct
\texttt{Cache-Control}, \texttt{ETag}, and \texttt{Vary} headers
(Chapters~\ref{ch:http_wire}, \ref{ch:idempotency}) translate directly into
origin shielding \cite{rfc9110}.

\textbf{Gateway.} The gateway is a policy enforcement point, not a junk
drawer. It offloads token verification (OAuth~2.0 bearer tokens with JWTs
\cite{rfc6749,rfc7519}, PKCE for public clients \cite{rfc7636}), enforces
rate tiers and emits IETF \texttt{RateLimit} headers
\cite{ietf_ratelimit_headers} (Chapter~\ref{ch:rate_limiting}), and stamps
every request with a correlation ID for tracing
(Chapter~\ref{ch:observability}). Business logic never lives here.

\textbf{Services.} Each service owns its profile: FastAPI REST for catalog
and orders with keyset pagination (Chapter~\ref{ch:pagination}), date-based
versioning (Chapter~\ref{ch:versioning}), RFC~9457 problem details
\cite{rfc9457} (Chapter~\ref{ch:problem_details}), and idempotency keys on
order creation (Chapter~\ref{ch:idempotency}). GraphQL fronts the read model
for the mobile app with depth limits and persisted queries
(Chapter~\ref{ch:graphql}). gRPC handles the internal inventory hot path
(Chapter~\ref{ch:grpc}). The dispatcher owns webhook signing, delivery logs,
and SSE fan-out (Chapter~\ref{ch:async_apis}).

\textbf{Data.} PostgreSQL is the system of record; every state change that
must be published is written \emph{in the same transaction} to an outbox
table, then relayed to the broker --- the transactional outbox pattern of
Chapter~\ref{ch:async_apis}, which eliminates dual-write anomalies
\cite{kleppmann2017ddia}. Redis carries rate-limit counters and short-lived
cache entries; it is never a source of truth.

\textbf{Observability.} OpenTelemetry instrumentation in every hop produces
traces joined by the gateway's correlation ID, RED/USE metrics, and
structured logs \cite{opentelemetrydocs}
(Chapter~\ref{ch:observability}). SLOs and error budgets follow the SRE
discipline of Chapter~\ref{ch:observability} \cite{beyer2016sre}.

\textbf{Governance.} OpenAPI and AsyncAPI specs are the source of truth,
linted in CI, rendered to the developer portal, and checked by the
conformance policy engine of Listing~3 (Section~\ref{sec:code-checker})
\cite{openapi31,asyncapi} (Chapter~\ref{ch:lifecycle}). Deprecations follow
the \texttt{Sunset} header convention \cite{rfc8594}.

\subsubsection{Request Lifecycle Walkthrough}

A partner's \texttt{POST /v1/orders} traverses the platform like this:

\begin{enumerate}
  \item \textbf{Edge:} TLS 1.3 handshake (resumed), WAF pass, cache bypass
        (POST is not cacheable \cite{rfc9110}).
  \item \textbf{Gateway:} JWT verified locally against cached JWKS
        \cite{rfc7519}; rate tier \texttt{partner-standard} admits the
        request, \texttt{RateLimit-Remaining: 97} stamped
        \cite{ietf_ratelimit_headers}; correlation ID \texttt{rid-9f2}
        assigned.
  \item \textbf{Service:} FastAPI validates the payload with Pydantic
        \cite{pydanticdocs,fastapidocs}; the \texttt{Idempotency-Key}
        deduplicates retries (Chapter~\ref{ch:idempotency}); a gRPC call to
        inventory reserves stock within the 50\,ms budget
        (Chapter~\ref{ch:grpc}); order row and outbox row commit in one
        transaction.
  \item \textbf{Async fan-out:} the outbox relay publishes
        \texttt{order.created} to the broker; the dispatcher signs and POSTs
        the webhook to the partner and pushes an SSE event to the operations
        dashboard (Chapter~\ref{ch:async_apis}).
  \item \textbf{Observe:} one trace, four spans, zero log lines containing
        credentials (Chapters~\ref{ch:observability},
        \ref{ch:api_security}).
\end{enumerate}

Every failure along this path has a named, rehearsed response: validation
fails with RFC~9457 422 \cite{rfc9457}; rate exhaustion returns 429 with
\texttt{Retry-After}; inventory timeout trips a circuit breaker
\cite{nygard2018releaseit} (Chapter~\ref{ch:microservices}); webhook delivery
retries with backoff and lands in the delivery log for replay.

\subsection{The Grand Capstone: Platform Specification}
\label{sec:grand-capstone}

This subsection is the full buildable specification of the platform in
Figure~\ref{fig:ref-arch}. It extends every per-chapter capstone into one
integrative build, and it is the [staff] project of
Section~\ref{sec:capstones}.

\subsubsection{Scope and Domains}

Build the Northwind Robotics API Platform as a set of deployable services
plus the governance toolchain that keeps them honest. Four domains:
\textbf{catalog} (products, categories, pricing), \textbf{orders} (placement,
status, partner notification), \textbf{inventory} (stock reservation,
restock), and \textbf{telemetry} (fleet readings, operator dashboards).
Synthetic data only; all services must run offline on one machine with
in-process fakes for the broker and payment gateway.

\subsubsection{API Surfaces per Protocol}

\begin{table}[htbp]
  \centering
  \small
  \begin{tabularx}{\textwidth}{l l X}
    \toprule
    \textbf{Surface} & \textbf{Profile} & \textbf{Contract} \\
    \midrule
    \texttt{/v1/products}, \texttt{/v1/orders} & REST + OpenAPI &
      keyset pagination, date version, RFC~9457 errors, idempotency keys \\
    \texttt{/graphql} & GraphQL &
      read model over catalog + telemetry; depth $\leq$ 5; persisted queries \\
    \texttt{inventory.v1.Inventory} & gRPC &
      \texttt{Reserve}/\texttt{Release}/\texttt{WatchStock} (server stream) \\
    \texttt{order.created}, \texttt{order.shipped} & webhooks &
      signed (HMAC), delivery log, replay endpoint, AsyncAPI spec \\
    \texttt{/v1/telemetry/stream} & SSE &
      operator dashboard feed, \texttt{Last-Event-ID} resume \\
    \bottomrule
  \end{tabularx}
  \caption{Grand Capstone API surfaces. Each surface cites its governing
  spec artifact.}
  \label{tab:capstone-surfaces}
\end{table}

\subsubsection{Non-Functional Requirements}

\begin{itemize}
  \item \textbf{SLOs:} catalog read p99 $<$ 200\,ms at 99.9\,\% monthly
        availability; order write p99 $<$ 500\,ms; inventory gRPC p99 $<$
        50\,ms; webhook delivery 99.5\,\% within 60\,s. Error budgets govern
        release pace \cite{beyer2016sre}.
  \item \textbf{Security posture:} OAuth~2.0 + JWT bearer at the gateway
        \cite{rfc6749,rfc7519}; API keys for partners; webhook signatures;
        OWASP API Top 10 checklist in CI \cite{owaspapisecurity}; no
        credential material in logs; TLS 1.3 only \cite{rfc8446}.
  \item \textbf{Rate tiers:} \texttt{public} 60/min, \texttt{partner-standard}
        600/min, \texttt{partner-premium} 6\,000/min; IETF headers on every
        response \cite{ietf_ratelimit_headers} (Chapter~\ref{ch:rate_limiting}).
  \item \textbf{Data integrity:} order + outbox atomicity; idempotent order
        placement; replayable webhook log retained 30 days
        \cite{kleppmann2017ddia}.
  \item \textbf{Operability:} OTel traces end to end \cite{opentelemetrydocs};
        twelve-factor configuration \cite{twelvefactor}; health endpoints;
        benchmark snapshot per release.
\end{itemize}

\subsubsection{Phased Delivery Plan}

\begin{table}[htbp]
  \centering
  \small
  \begin{tabularx}{\textwidth}{l l X}
    \toprule
    \textbf{Milestone} & \textbf{Theme} & \textbf{Exit criteria} \\
    \midrule
    M1 (weeks 1--4) & REST core &
      catalog + orders REST surface green on contract tests; RFC~9457
      everywhere; keyset pagination; idempotency replay test passing. \\
    M2 (weeks 5--8) & Read model + internal RPC &
      GraphQL read model live with depth limiting; inventory gRPC service
      with streaming \texttt{WatchStock}; gateway auth + rate tiers. \\
    M3 (weeks 9--12) & Async platform &
      outbox relay, signed webhooks with delivery log and replay, SSE
      telemetry stream; AsyncAPI specs linted in CI. \\
    M4 (weeks 13--16) & Production hardening &
      OTel traces across all surfaces; SLO dashboards; conformance checker
      gating merges; load and failure-mode test evidence; portal published. \\
    \bottomrule
  \end{tabularx}
  \caption{Grand Capstone four-milestone delivery plan.}
  \label{tab:milestones}
\end{table}

\subsubsection{Acceptance Criteria and Evaluation Rubric}

Acceptance is binary and evidence-based: every surface in
Table~\ref{tab:capstone-surfaces} responds per its spec in an offline
in-process test run; the conformance checker passes on the manifests folder;
the SLO dashboard shows green against a recorded load test; a chaos drill
(kill inventory mid-order) produces the designed circuit-breaker behavior
\cite{nygard2018releaseit}.

The rubric scores five dimensions at 0--4 each: \textbf{correctness}
(contract conformance), \textbf{resilience} (failure-mode behavior),
\textbf{operability} (telemetry, runbooks), \textbf{governance} (specs,
lint, portal), and \textbf{craft} (code quality, typing, tests).
A passing Grand Capstone scores $\geq$ 16 with no dimension below 3.

\subsection{Cross-Cutting Synthesis}
\label{sec:synthesis}

\subsubsection{The Failure Taxonomy, Revisited}

Chapter by chapter, this book has built a taxonomy of how APIs fail.
Table~\ref{tab:failure-taxonomy} consolidates it: for each failure class,
where it is taught, how you detect it, and the designed mitigation.

\begin{table}[htbp]
  \centering
  \small
  \begin{tabularx}{\textwidth}{l l X X}
    \toprule
    \textbf{Failure class} & \textbf{Chapter} & \textbf{Detection} & \textbf{Mitigation} \\
    \midrule
    Protocol misuse & \ref{ch:http_wire} & semantic lint, packet capture & teach HTTP, not frameworks \cite{rfc9110} \\
    Client brittleness & \ref{ch:consuming_python} & contract tests & retries with backoff, timeouts \cite{httpxdocs} \\
    Ambiguous errors & \ref{ch:problem_details} & error-shape lint & RFC~9457 problem details \cite{rfc9457} \\
    Page drift & \ref{ch:pagination} & duplicate/skipped rows & keyset cursors \\
    Breaking change & \ref{ch:versioning} & spec diff in CI & date versions, \texttt{Sunset} \cite{rfc8594} \\
    Duplicate side effects & \ref{ch:idempotency} & replay tests & idempotency keys \\
    Lost events & \ref{ch:async_apis} & outbox lag, delivery log & transactional outbox \cite{kleppmann2017ddia} \\
    N+1 resolvers & \ref{ch:graphql} & query tracing & dataloaders, depth limits \\
    Tight RPC coupling & \ref{ch:grpc} & stub fan-out audit & contract-first, versioned protos \\
    Quota abuse & \ref{ch:rate_limiting} & 429 rate, tier stats & tiers + IETF headers \cite{ietf_ratelimit_headers} \\
    Cascading failure & \ref{ch:microservices} & latency waterfall & timeouts, breakers, bulkheads \cite{nygard2018releaseit} \\
    Credential leaks & \ref{ch:api_security} & log scans, secret scans & OWASP discipline \cite{owaspapisecurity} \\
    Snowflake deploys & \ref{ch:deployment} & drift detection & twelve-factor config \cite{twelvefactor} \\
    Blind operation & \ref{ch:observability} & missing spans & OTel everywhere \cite{opentelemetrydocs} \\
    Untested contracts & \ref{ch:testing} & consumer complaints & consumer-driven contracts \cite{pactdocs} \\
    Spec rot & \ref{ch:lifecycle} & portal/spec drift & executable governance \cite{openapi31,asyncapi} \\
    \bottomrule
  \end{tabularx}
  \caption{The consolidated failure taxonomy: one row per chapter's core
  enemy.}
  \label{tab:failure-taxonomy}
\end{table}

\subsubsection{One Page per Module}

\begin{table}[htbp]
  \centering
  \small
  \begin{tabularx}{\textwidth}{l X X}
    \toprule
    \multicolumn{3}{l}{\textbf{Module I --- Foundations (Ch.~\ref{ch:introduction}--\ref{ch:client_security})}} \\
    \midrule
    \textbf{Chapter} & \textbf{Core idea} & \textbf{Artifact you must produce} \\
    \midrule
    \ref{ch:introduction} & APIs as products and contracts & platform vision + consumer map \\
    \ref{ch:http_wire} & messages, semantics, connections \cite{rfc9110,rfc9112,rfc9113} & annotated request/response traces \\
    \ref{ch:consuming_python} & disciplined clients & retrying, timeout-bound client \cite{httpxdocs,requestsdocs} \\
    \ref{ch:rest_style} & resources, uniformity, maturity \cite{fielding2000rest,richardson2013restful} & resource model + OpenAPI draft \\
    \ref{ch:serialization} & contracts over bytes \cite{rfc8259,jsonschemadocs} & versioned schemas \\
    \ref{ch:client_security} & secrets, TLS, token handling \cite{rfc8446,rfc7636} & secure client checklist \\
    \bottomrule
  \end{tabularx}
  \caption{Module I synthesis: speak the protocol fluently before building.}
  \label{tab:module1}
\end{table}

\begin{table}[htbp]
  \centering
  \small
  \begin{tabularx}{\textwidth}{l X X}
    \toprule
    \multicolumn{3}{l}{\textbf{Module II --- Building APIs (Ch.~\ref{ch:fastapi}--\ref{ch:idempotency})}} \\
    \midrule
    \textbf{Chapter} & \textbf{Core idea} & \textbf{Artifact you must produce} \\
    \midrule
    \ref{ch:fastapi} & typed, fast service scaffolding \cite{fastapidocs,pydanticdocs} & working FastAPI service \\
    \ref{ch:problem_details} & errors as a designed interface \cite{rfc9457} & problem-details catalog \\
    \ref{ch:pagination} & stable traversal of changing data \cite{rfc8288} & keyset-paginated collections \\
    \ref{ch:versioning} & evolve without breaking & date versioning + \texttt{Sunset} policy \cite{rfc8594} \\
    \ref{ch:idempotency} & retries must be safe & idempotency-key middleware \\
    \bottomrule
  \end{tabularx}
  \caption{Module II synthesis: correctness is a designed property.}
  \label{tab:module2}
\end{table}

\begin{table}[htbp]
  \centering
  \small
  \begin{tabularx}{\textwidth}{l X X}
    \toprule
    \multicolumn{3}{l}{\textbf{Module III --- Protocols \& Scale (Ch.~\ref{ch:async_apis}--\ref{ch:microservices})}} \\
    \midrule
    \textbf{Chapter} & \textbf{Core idea} & \textbf{Artifact you must produce} \\
    \midrule
    \ref{ch:async_apis} & decouple with events and push \cite{asyncapi} & outbox + webhook pipeline \\
    \ref{ch:graphql} & client-driven selection \cite{graphqlspec} & governed schema + persisted queries \\
    \ref{ch:grpc} & typed, streaming RPC \cite{grpcdocs,protobufdocs} & proto contracts + stubs \\
    \ref{ch:rate_limiting} & protect capacity, signal policy \cite{ietf_ratelimit_headers} & tiered limiter + headers \\
    \ref{ch:microservices} & resilience is engineered \cite{nygard2018releaseit,newman2021microservices} & breaker/timeout budget map \\
    \bottomrule
  \end{tabularx}
  \caption{Module III synthesis: pick profiles per constraint, scale
  deliberately.}
  \label{tab:module3}
\end{table}

\begin{table}[htbp]
  \centering
  \small
  \begin{tabularx}{\textwidth}{l X X}
    \toprule
    \multicolumn{3}{l}{\textbf{Module IV --- Production \& Governance (Ch.~\ref{ch:api_security}--\ref{ch:lifecycle})}} \\
    \midrule
    \textbf{Chapter} & \textbf{Core idea} & \textbf{Artifact you must produce} \\
    \midrule
    \ref{ch:api_security} & defense in depth for APIs \cite{owaspapisecurity,rfc6749,rfc7519} & threat model + auth design \\
    \ref{ch:deployment} & reproducible, portable services \cite{twelvefactor} & container + CI pipeline \\
    \ref{ch:observability} & you cannot run what you cannot see \cite{opentelemetrydocs,beyer2016sre} & traces, SLOs, runbooks \\
    \ref{ch:testing} & contracts are the test surface \cite{pactdocs} & contract + chaos suites \\
    \ref{ch:lifecycle} & governance as code \cite{openapi31,gehl2021apidesign} & linted specs, portal, policy CI \\
    \bottomrule
  \end{tabularx}
  \caption{Module IV synthesis: production is a discipline, not a
  destination.}
  \label{tab:module4}
\end{table}

\subsection{Emerging Frontiers, Honestly Assessed}
\label{sec:frontiers}

A capstone chapter owes you more than enthusiasm: it owes you an
adopt/trial/wait posture for each frontier, with the maturity caveats that
vendor keynotes omit.

\subsubsection{HTTP/3 and QUIC}

HTTP/3 \cite{rfc9114} reimplements HTTP semantics over QUIC \cite{rfc9000},
which replaces TCP + TLS 1.3 \cite{rfc8446} with a UDP-based transport that
integrates the handshake, provides per-stream loss recovery (eliminating
TCP's head-of-line blocking that HTTP/2 \cite{rfc9113,rfc7541} merely
mitigates), and supports connection migration across network changes ---
a direct win for mobile clients moving between radios. \textbf{Posture:
trial.} Browsers and major CDNs already negotiate HTTP/3 via
\texttt{Alt-Svc}, so public REST and SSE surfaces inherit it for free at the
edge; but QUIC's CPU cost at the server, UDP filtering in enterprise
networks, and thinner debugging tooling mean you should measure before
mandating it. gRPC stays on HTTP/2 in most stacks today.

\subsubsection{APIs Consumed by AI Agents (MCP)}

The newest consumer class is not a human developer or a mobile app but an AI
agent invoking tools. The Model Context Protocol (MCP) standardizes how hosts
discover and call such tools: a server publishes \emph{tool schemas} ---
typed names, descriptions, and JSON Schema argument contracts
\cite{jsonschemadocs} --- and the agent selects and invokes them at runtime.
The implication for API design is profound: \textbf{your schema is now read
by a model, not just by humans}. Descriptions, field naming, error shapes,
and idempotency semantics become prompt-adjacent surface area. Tool schemas
should be derived from the same OpenAPI/JSON Schema source of truth
\cite{openapi31} rather than hand-maintained, or the two contracts will
drift --- exactly the spec-rot failure of Chapter~\ref{ch:lifecycle} in a new
costume. RFC~9457 problem details \cite{rfc9457} turn out to be
agent-friendly errors: structured, typed, and self-describing.
\textbf{Posture: trial} for internal tool ecosystems; \textbf{wait} before
exposing MCP surfaces publicly until authorization and auditing patterns
settle.

\subsubsection{AI-Assisted API Design and Review}

LLM assistants now draft OpenAPI documents, generate SDKs, and review spec
diffs competently. The durable value is not generation --- it is
\emph{review at scale}: an assistant can flag missing error responses,
inconsistent pagination, or non-idiomatic naming across hundreds of
endpoints. The durable risk is confident invention of plausible-but-wrong
contracts. The professional pattern: AI proposes, executable policy (the
conformance checker of Listing~3 (Section~\ref{sec:code-checker})) disposes. \textbf{Posture:
adopt} for drafting and review, always behind deterministic CI gates.

\subsubsection{Edge-Native APIs}

Edge runtimes (worker isolates, regional endpoints) move request handling
closer to the consumer, cutting TLS and TTFB latency and enabling regional
data affinity for compliance. The architectural constraint is state: edge
functions shine for read-heavy, cache-friendly, or geographically partitioned
workloads and hurt when every request must reach a centralized strongly
consistent store. Regional affinity --- routing a tenant's writes to their
home region --- is the pragmatic middle. \textbf{Posture: trial} for
read-dominated surfaces (catalog!); \textbf{wait} for write-heavy cores.

\subsubsection{Post-Quantum TLS Readiness}

Harvest-now-decrypt-later attacks make cryptographic agility a present-tense
concern. The ecosystem's answer is hybrid key exchange in TLS~1.3
\cite{rfc8446} (classical + post-quantum KEM in one handshake), already
rolling out across browsers and CDNs. For API platforms the work is
unglamorous: inventory where TLS is terminated, ensure libraries track
current hybrids, watch handshake-size regressions that can break QUIC
amplification limits, and keep certificate automation healthy so algorithm
rotation is an upgrade, not a migration. \textbf{Posture: adopt} at the edge
(inherited from your CDN); \textbf{trial} for east--west service mesh TLS.

\subsubsection{API Security Automation}

Attack surface discovery, schema-aware fuzzing, and runtime anomaly detection
are converging into continuous API security testing keyed off your OpenAPI
specs \cite{openapi31,owaspapisecurity}. The spec-first discipline of
Chapter~\ref{ch:lifecycle} is the prerequisite: tools can only fuzz what they
can see. \textbf{Posture: adopt} spec-driven scanning in CI;
\textbf{trial} runtime behavioral detection.

\begin{table}[htbp]
  \centering
  \small
  \begin{tabularx}{\textwidth}{l l X}
    \toprule
    \textbf{Frontier} & \textbf{Posture} & \textbf{Prerequisite before adoption} \\
    \midrule
    HTTP/3 / QUIC \cite{rfc9114,rfc9000} & trial & edge support inherited; measure server CPU \\
    MCP tool schemas & trial (internal) & spec-first JSON Schema source of truth \cite{jsonschemadocs} \\
    AI-assisted design/review & adopt & deterministic lint gates behind every suggestion \\
    Edge-native APIs & trial & read-heavy or regionally partitionable data \\
    Post-quantum TLS & adopt (edge) & TLS inventory, automation, current libraries \cite{rfc8446} \\
    Security automation & adopt (CI) & complete, linted OpenAPI coverage \cite{owaspapisecurity} \\
    \bottomrule
  \end{tabularx}
  \caption{Frontier adoption postures. Nothing here exempts you from the
  operational maturity of Module IV.}
  \label{tab:frontiers}
\end{table}

\begin{warningbox}
\textbf{Frontiers amplify your existing maturity.} HTTP/3 does not fix a
missing caching strategy; MCP does not fix an unversioned schema; edge
compute does not fix unbounded fan-out. Adopt frontiers \emph{after} the
boring disciplines of Chapters~\ref{ch:api_security}--\ref{ch:lifecycle} are
routine, or you will simply fail faster, in more places.
\end{warningbox}

\subsection{Career Pathways and the Mastery Checklist}
\label{sec:careers}

API engineering is a T-shaped career: broad literacy across protocols and
operations, deep expertise that grows with each level.
Figure~\ref{fig:pathway} maps the levels onto the book's modules;
Table~\ref{tab:competencies} makes the progression checkable.

\begin{figure}[htbp]
\centering
\begin{tikzpicture}[
  font=\small,
  >=stealth,
  node distance=6mm,
  lvl/.style={flowstep, fill=blue!5, align=center, minimum width=3.1cm, minimum height=1.5cm},
  lbl/.style={modcap}
]
\node[lvl] (junior) {\textbf{Junior}\\[2pt] consume + build\\basic REST\\Ch.~\ref{ch:introduction}--\ref{ch:fastapi}};
\node[lvl, right=of junior] (mid) {\textbf{Mid}\\[2pt] correctness:\\errors, pagination,\\versions, idempotency\\Ch.~\ref{ch:problem_details}--\ref{ch:idempotency}};
\node[lvl, right=of mid] (senior) {\textbf{Senior}\\[2pt] multi-protocol,\\resilience, scale\\Ch.~\ref{ch:async_apis}--\ref{ch:microservices}};
\node[lvl, right=of senior] (staff) {\textbf{Staff}\\[2pt] platform, security,\\governance, strategy\\Ch.~\ref{ch:api_security}--\ref{ch:lifecycle}, this ch.};
\draw[->, thick] (junior) -- (mid);
\draw[->, thick] (mid) -- (senior);
\draw[->, thick] (senior) -- (staff);
\node[lbl, above=2mm of mid] {increasing scope: endpoint $\rightarrow$ service $\rightarrow$ system $\rightarrow$ organization};
\end{tikzpicture}
\caption{The API engineer mastery pathway. Levels are scopes of
responsibility, not years of tenure.}
\label{fig:pathway}
\end{figure}

\begin{table}[htbp]
  \centering
  \small
  \begin{tabularx}{\textwidth}{l X}
    \toprule
    \textbf{Level} & \textbf{You can be trusted to\ldots} \\
    \midrule
    Junior &
      read a packet capture; consume any documented REST API with retries
      and timeouts; build a typed FastAPI CRUD service; keep secrets out of
      code and logs. \\
    Mid &
      design RFC~9457 error catalogs; implement keyset pagination and date
      versioning; make every mutation idempotent; write contract tests;
      explain your API's caching headers to a CDN engineer. \\
    Senior &
      choose protocols from constraints and defend the choice; design
      outbox-based event pipelines; operate rate tiers; break production
      safely (timeouts, breakers, bulkheads) and prove it with chaos
      drills \cite{nygard2018releaseit}. \\
    Staff &
      own a multi-protocol platform strategy; set SLOs and error budgets
      \cite{beyer2016sre}; run governance as code; evaluate frontiers with
      adopt/trial/wait rigor; grow seniors through design review
      \cite{gehl2021apidesign}. \\
    \bottomrule
  \end{tabularx}
  \caption{The mastery checklist, by level.}
  \label{tab:competencies}
\end{table}

\textbf{Staying current.} The durable sources are the specification bodies
and their errata, not social media: the IETF HTTP and QUIC working groups
(RFCs 9110--9114 \cite{rfc9110,rfc9114}), the OpenAPI Initiative
\cite{openapi31}, the AsyncAPI Initiative \cite{asyncapi}, the GraphQL
Foundation \cite{graphqlspec}, the CNCF (gRPC \cite{grpcdocs}, OpenTelemetry
\cite{opentelemetrydocs}), and OWASP \cite{owaspapisecurity}. Conferences
worth your time: API World, APIDays, KubeCon/CloudNativeCon, QCon, and SREcon.
Communities: vendor-neutral spec discussion lists and the conference talk
archives. Budget one hour a week; reread foundational texts
\cite{fielding2000rest,kleppmann2017ddia,newman2021microservices} yearly ---
they reward rereading as your scope grows.

\section{Production-Ready Code Implementation}
\label{sec:code}

Three integrative, complete programs implement this chapter's ideas. All are
offline-first (no network, no credentials), deterministic (no wall-clock or
randomness on decision paths), fully typed, and shipped with executable
tests. Setup (PowerShell, run once per machine):

\begin{minted}{powershell}
cd "code\ch21"
python -m venv .venv
.\.venv\Scripts\Activate.ps1
pip install fastapi "strawberry-graphql[fastapi]" pydantic pyyaml httpx pytest
python -m pytest -q
\end{minted}

\subsection{Listing 1 --- The Protocol-Selection Advisor}
\label{sec:code-advisor}

The advisor encodes the Master Decision Card of
Section~\ref{sec:decision-card} as a scored questionnaire. Each criterion
emits \emph{signals}: weighted, human-readable point contributions to one or
more profiles. The score for profile $p$ is simply

\begin{equation}
\mathrm{score}(p) = \sum_{s \in S,\; s.\mathrm{profile} = p} s.\mathrm{points}
\end{equation}

and the recommendation is the argmax with a deterministic tie-break by
profile declaration order. Because every signal carries its rationale, the
recommendation ships with an \emph{auditable trace} --- you can show a
reviewer exactly which constraints moved which profile, which is how the
framework stays teachable instead of becoming an oracle.

\begin{minted}{python}
"""Protocol-selection advisor for the Master Decision Card.

Encodes the Chapter 21 decision framework as a scored questionnaire:
weighted criteria, deterministic arithmetic, and a full rationale trace
so the recommendation can be audited and taught from.

Offline, deterministic, standard library only. Python 3.11+.
"""

from __future__ import annotations

from dataclasses import dataclass, field
from enum import Enum


class Profile(str, Enum):
    """API interaction profiles compared by the advisor."""

    REST = "REST"
    GRAPHQL = "GraphQL"
    GRPC = "gRPC"
    EVENTS = "Event-driven (webhooks/broker)"
    STREAM = "WebSockets / SSE"


class Consumer(str, Enum):
    PUBLIC = "public-third-party"
    MOBILE = "mobile-spa"
    INTERNAL = "internal-service"
    MACHINE = "machine-telemetry"


class Payload(str, Enum):
    FIXED = "fixed-resources"
    GRAPH = "nested-graph"
    HIGH_VOLUME = "high-volume-streams"


class Latency(str, Enum):
    RELAXED = "relaxed"      # p99 >= 200 ms is fine
    TIGHT = "tight"          # interactive, p99 < 200 ms
    EXTREME = "extreme"      # p99 < 50 ms service-to-service


class Delivery(str, Enum):
    SYNC = "sync-request-response"
    PUSH = "server-push-realtime"
    ASYNC = "async-fire-and-forget"


class Team(str, Enum):
    EXPLORATORY = "exploratory"
    DISCIPLINED = "disciplined"


@dataclass(frozen=True)
class Requirements:
    """Answers to the six-question selection questionnaire."""

    consumer: Consumer
    payload: Payload
    latency: Latency
    delivery: Delivery
    http_caching: bool
    team: Team


@dataclass(frozen=True)
class Signal:
    """One weighted contribution to a profile's score."""

    profile: Profile
    points: int
    criterion: str
    rationale: str


@dataclass(frozen=True)
class Recommendation:
    """Advisor output: ranked scores plus an auditable trace."""

    winner: Profile
    scores: dict[Profile, int]
    signals: tuple[Signal, ...] = field(default_factory=tuple)

    def rationale(self, profile: Profile | None = None) -> list[str]:
        """Return the trace lines for one profile (default: the winner)."""
        target = profile or self.winner
        return [
            f"[{s.points:+d}] {s.criterion}: {s.rationale}"
            for s in self.signals
            if s.profile == target
        ]


def _score(req: Requirements) -> tuple[Signal, ...]:
    """Apply every rule of the Master Decision Card to the answers."""
    signals: list[Signal] = []

    def add(profile: Profile, points: int, criterion: str, rationale: str) -> None:
        signals.append(Signal(profile, points, criterion, rationale))

    # --- Criterion 1: consumer type -------------------------------------
    if req.consumer is Consumer.PUBLIC:
        add(Profile.REST, 3, "consumer",
            "public third parties expect plain HTTPS + JSON they can curl")
        add(Profile.EVENTS, 2, "consumer",
            "webhooks are the standard public async contract")
        add(Profile.GRAPHQL, 1, "consumer",
            "powerful for public developers but raises the support burden")
        add(Profile.GRPC, -3, "consumer",
            "browser clients cannot speak raw HTTP/2 gRPC without a proxy")
        add(Profile.STREAM, 1, "consumer",
            "SSE/WebSocket endpoints are feasible but operationally niche")
    elif req.consumer is Consumer.MOBILE:
        add(Profile.GRAPHQL, 3, "consumer",
            "clients pick exactly the fields they render; no over-fetching")
        add(Profile.REST, 2, "consumer",
            "aggregating BFF-style REST resources also solve over-fetching")
        add(Profile.STREAM, 1, "consumer",
            "push channels pair well with mobile UX")
        add(Profile.GRPC, -2, "consumer",
            "gRPC-Web adds a proxy hop that mobile teams rarely want")
    elif req.consumer is Consumer.INTERNAL:
        add(Profile.GRPC, 3, "consumer",
            "strongly typed, streaming, low-overhead service-to-service RPC")
        add(Profile.REST, 1, "consumer",
            "REST remains the default when tooling maturity is low")
        add(Profile.EVENTS, 1, "consumer",
            "async integration decouples internal services")
        add(Profile.GRAPHQL, 1, "consumer",
            "federated graphs can front many internal services")
    else:  # MACHINE telemetry
        add(Profile.EVENTS, 3, "consumer",
            "telemetry is naturally asynchronous and bursty")
        add(Profile.GRPC, 2, "consumer",
            "bidirectional streaming fits high-rate device channels")
        add(Profile.STREAM, 1, "consumer",
            "persistent sockets work where HTTP/2 is unavailable")
        add(Profile.REST, -1, "consumer",
            "per-reading REST calls waste headers and connections")

    # --- Criterion 2: payload shape -------------------------------------
    if req.payload is Payload.GRAPH:
        add(Profile.GRAPHQL, 3, "payload",
            "nested graph selection sets are GraphQL's home turf")
        add(Profile.REST, -1, "payload",
            "deep graphs force N+1 endpoint hops in REST")
    elif req.payload is Payload.HIGH_VOLUME:
        add(Profile.GRPC, 2, "payload",
            "compact binary frames and HTTP/2 streams suit high volume")
        add(Profile.EVENTS, 2, "payload",
            "brokered batches absorb volume spikes")
        add(Profile.REST, -1, "payload",
            "textual JSON per request is the most expensive option")
    else:
        add(Profile.REST, 2, "payload",
            "fixed resource shapes map one-to-one onto REST resources")

    # --- Criterion 3: latency budget ------------------------------------
    if req.latency is Latency.EXTREME:
        add(Profile.GRPC, 3, "latency",
            "HTTP/2 multiplexing plus Protobuf minimize per-call overhead")
        add(Profile.GRAPHQL, -1, "latency",
            "query planning and JSON add milliseconds per request")
        add(Profile.REST, 0, "latency",
            "acceptable with keep-alive and HTTP/2, but not minimal")
    elif req.latency is Latency.TIGHT:
        add(Profile.GRPC, 1, "latency", "comfortable within tight budgets")
        add(Profile.REST, 1, "latency", "comfortable within tight budgets")
    else:
        add(Profile.EVENTS, 2, "latency",
            "async delivery tolerates relaxed latency by design")

    # --- Criterion 4: delivery model (hard structural signals) ----------
    if req.delivery is Delivery.PUSH:
        add(Profile.STREAM, 4, "delivery",
            "server push requires SSE or WebSockets; unary APIs cannot push")
        add(Profile.GRPC, 1, "delivery",
            "server-streaming RPC works for non-browser clients")
        add(Profile.REST, -3, "delivery",
            "unary REST has no server-push channel")
        add(Profile.GRAPHQL, -2, "delivery",
            "GraphQL subscriptions need a socket transport anyway")
    elif req.delivery is Delivery.ASYNC:
        add(Profile.EVENTS, 4, "delivery",
            "fire-and-forget delivery is the defining event-driven case")
    else:
        add(Profile.REST, 1, "delivery",
            "synchronous request-response is REST's native shape")
        add(Profile.GRPC, 1, "delivery",
            "unary RPC is also synchronous request-response")

    # --- Criterion 5: HTTP caching --------------------------------------
    if req.http_caching:
        add(Profile.REST, 3, "caching",
            "only REST leverages shared/CDN caching of GET responses")
        add(Profile.GRAPHQL, -2, "caching",
            "POST-based queries bypass standard HTTP caches")
        add(Profile.GRPC, -2, "caching",
            "gRPC calls are not addressable cacheable resources")
    else:
        add(Profile.GRAPHQL, 1, "caching",
            "no penalty when edge caching is not required")

    # --- Criterion 6: team maturity -------------------------------------
    if req.team is Team.EXPLORATORY:
        add(Profile.REST, 2, "team",
            "REST has the shallowest learning curve and richest tooling")
        add(Profile.GRPC, -1, "team",
            "Protobuf toolchains and reflection quirks slow new teams")
        add(Profile.GRAPHQL, -1, "team",
            "schema governance and N+1 discipline require experience")
    else:
        add(Profile.GRAPHQL, 1, "team",
            "disciplined teams can exploit schema-first workflows")
        add(Profile.GRPC, 1, "team",
            "disciplined teams can run contract-first RPC safely")

    return tuple(signals)


def recommend(req: Requirements) -> Recommendation:
    """Score all profiles and return the winner with a rationale trace.

    Ties break deterministically by profile enum order, so the same
    answers always produce the same recommendation.
    """
    signals = _score(req)
    scores = {profile: 0 for profile in Profile}
    for signal in signals:
        scores[signal.profile] += signal.points
    winner = max(Profile, key=lambda p: (scores[p], -list(Profile).index(p)))
    return Recommendation(winner=winner, scores=scores, signals=signals)


# --- Canonical scenarios from the chapter's worked examples ---------------

PUBLIC_SAAS = Requirements(
    consumer=Consumer.PUBLIC,
    payload=Payload.FIXED,
    latency=Latency.RELAXED,
    delivery=Delivery.SYNC,
    http_caching=True,
    team=Team.DISCIPLINED,
)

INTERNAL_MESH = Requirements(
    consumer=Consumer.INTERNAL,
    payload=Payload.FIXED,
    latency=Latency.EXTREME,
    delivery=Delivery.SYNC,
    http_caching=False,
    team=Team.DISCIPLINED,
)

MOBILE_BACKEND = Requirements(
    consumer=Consumer.MOBILE,
    payload=Payload.GRAPH,
    latency=Latency.TIGHT,
    delivery=Delivery.SYNC,
    http_caching=False,
    team=Team.DISCIPLINED,
)

IOT_TELEMETRY = Requirements(
    consumer=Consumer.MACHINE,
    payload=Payload.HIGH_VOLUME,
    latency=Latency.RELAXED,
    delivery=Delivery.ASYNC,
    http_caching=False,
    team=Team.DISCIPLINED,
)

PARTNER_NOTIFICATIONS = Requirements(
    consumer=Consumer.PUBLIC,
    payload=Payload.FIXED,
    latency=Latency.RELAXED,
    delivery=Delivery.ASYNC,
    http_caching=False,
    team=Team.EXPLORATORY,
)


if __name__ == "__main__":
    scenarios = {
        "Public SaaS API": PUBLIC_SAAS,
        "Internal microservice mesh": INTERNAL_MESH,
        "Mobile app backend": MOBILE_BACKEND,
        "IoT telemetry": IOT_TELEMETRY,
        "Partner notifications": PARTNER_NOTIFICATIONS,
    }
    for name, reqs in scenarios.items():
        rec = recommend(reqs)
        print(f"{name}: {rec.winner.value}")
        for line in rec.rationale():
            print(f"    {line}")
\end{minted}

The unit tests lock the five canonical worked scenarios to their expected
recommendations and assert the framework invariants: determinism, score
arithmetic (scores equal the sum of their signals), and trace integrity.

\begin{minted}{python}
"""Unit tests for the protocol-selection advisor.

Covers the five canonical scenarios from the chapter plus framework
invariants (determinism, scoring arithmetic, rationale trace).
"""

from __future__ import annotations

from protocol_advisor import (
    IOT_TELEMETRY,
    INTERNAL_MESH,
    MOBILE_BACKEND,
    PARTNER_NOTIFICATIONS,
    PUBLIC_SAAS,
    Consumer,
    Delivery,
    Latency,
    Payload,
    Profile,
    Requirements,
    Team,
    recommend,
)


def test_public_saas_recommends_rest() -> None:
    assert recommend(PUBLIC_SAAS).winner is Profile.REST


def test_internal_mesh_recommends_grpc() -> None:
    assert recommend(INTERNAL_MESH).winner is Profile.GRPC


def test_mobile_backend_recommends_graphql() -> None:
    assert recommend(MOBILE_BACKEND).winner is Profile.GRAPHQL


def test_iot_telemetry_recommends_events() -> None:
    assert recommend(IOT_TELEMETRY).winner is Profile.EVENTS


def test_partner_notifications_recommend_events() -> None:
    assert recommend(PARTNER_NOTIFICATIONS).winner is Profile.EVENTS


def test_push_delivery_selects_streams() -> None:
    reqs = Requirements(
        consumer=Consumer.PUBLIC,
        payload=Payload.FIXED,
        latency=Latency.TIGHT,
        delivery=Delivery.PUSH,
        http_caching=False,
        team=Team.DISCIPLINED,
    )
    assert recommend(reqs).winner is Profile.STREAM


def test_recommendation_is_deterministic() -> None:
    first = recommend(PUBLIC_SAAS)
    second = recommend(PUBLIC_SAAS)
    assert first == second


def test_scores_equal_sum_of_signals() -> None:
    rec = recommend(MOBILE_BACKEND)
    for profile in Profile:
        total = sum(s.points for s in rec.signals if s.profile is profile)
        assert rec.scores[profile] == total


def test_winner_has_strictly_highest_score() -> None:
    rec = recommend(IOT_TELEMETRY)
    best = max(rec.scores.values())
    assert rec.scores[rec.winner] == best
    assert list(rec.scores.values()).count(best) >= 1


def test_rationale_trace_references_winner_only() -> None:
    rec = recommend(INTERNAL_MESH)
    trace = rec.rationale()
    assert trace, "winner must have at least one rationale line"
    for line in trace:
        assert ":" in line and line.startswith("[")
\end{minted}

Running the demo driver prints the recommendation and the rationale trace
for each canonical scenario:

\begin{minted}{text}
Public SaaS API: REST
    [+3] consumer: public third parties expect plain HTTPS + JSON they can curl
    [+2] payload: fixed resource shapes map one-to-one onto REST resources
    [+1] delivery: synchronous request-response is REST's native shape
    [+3] caching: only REST leverages shared/CDN caching of GET responses
Internal microservice mesh: gRPC
    [+3] consumer: strongly typed, streaming, low-overhead service-to-service RPC
    [+3] latency: HTTP/2 multiplexing plus Protobuf minimize per-call overhead
    [+1] delivery: unary RPC is also synchronous request-response
    [+1] team: disciplined teams can run contract-first RPC safely
Mobile app backend: GraphQL
    [+3] consumer: clients pick exactly the fields they render; no over-fetching
    [+3] payload: nested graph selection sets are GraphQL's home turf
    [+1] caching: no penalty when edge caching is not required
    [+1] team: disciplined teams can exploit schema-first workflows
IoT telemetry: Event-driven (webhooks/broker)
    [+3] consumer: telemetry is naturally asynchronous and bursty
    [+2] payload: brokered batches absorb volume spikes
    [+2] latency: async delivery tolerates relaxed latency by design
    [+4] delivery: fire-and-forget delivery is the defining event-driven case
Partner notifications: Event-driven (webhooks/broker)
    [+2] consumer: webhooks are the standard public async contract
    [+2] latency: async delivery tolerates relaxed latency by design
    [+4] delivery: fire-and-forget delivery is the defining event-driven case
\end{minted}

\subsection{Listing 2 --- The Northwind Platform Slice}
\label{sec:code-platform}

One FastAPI application \cite{fastapidocs,pydanticdocs} assembling the
book's patterns into a single deployable unit: a keyset-paginated REST
catalog (Chapters~\ref{ch:fastapi}, \ref{ch:pagination}), a GraphQL read
model joining catalog rows to inventory availability
(Chapter~\ref{ch:graphql}), an SSE telemetry stream
(Chapter~\ref{ch:async_apis}), RFC~9457 problem-details errors from every
failure path \cite{rfc9457} (Chapter~\ref{ch:problem_details}), idempotency
keys on order creation (Chapter~\ref{ch:idempotency}), and IETF
\texttt{RateLimit} headers from a fixed-window limiter
\cite{ietf_ratelimit_headers} (Chapter~\ref{ch:rate_limiting}). State is
injected through the application factory so every test gets an isolated
platform --- the deterministic-test discipline of Chapter~\ref{ch:testing}.

\begin{minted}{python}
"""Northwind Robotics API Platform -- integrative slice (Chapter 21).

One FastAPI application assembling the book's core patterns into a
single deployable unit:

* REST catalog with keyset pagination          (Ch. 6, Ch. 8)
* GraphQL read model over the same domain data (Ch. 12)
* SSE telemetry stream                         (Ch. 11)
* RFC 9457 problem-details errors              (Ch. 7)
* Idempotency keys on order creation           (Ch. 10)
* IETF rate-limit response headers             (Ch. 14)

Offline, deterministic, synthetic data only. Python 3.11+.
"""

from __future__ import annotations

import base64
import threading
from collections.abc import AsyncIterator
from dataclasses import dataclass, field

import strawberry
import uvicorn
from fastapi import FastAPI, Query, Request
from fastapi.exceptions import RequestValidationError
from fastapi.responses import JSONResponse, StreamingResponse
from pydantic import BaseModel, Field
from strawberry.fastapi import GraphQLRouter

# --- Domain data (synthetic Northwind Robotics universe) -------------------

PRODUCTS: list[dict[str, object]] = [
    {
        "id": 1,
        "sku": "NWR-ARM-6X",
        "name": "Northwind 6-Axis Arm",
        "category": "manipulators",
        "unit_price_cents": 12_500_00,
    },
    {
        "id": 2,
        "sku": "NWR-LDR-360",
        "name": "LiDAR 360 Scanner",
        "category": "sensors",
        "unit_price_cents": 3_499_00,
    },
    {
        "id": 3,
        "sku": "NWR-DRV-BLDC",
        "name": "BLDC Motor Driver",
        "category": "actuators",
        "unit_price_cents": 189_99,
    },
]

AVAILABILITY: dict[str, int] = {
    "NWR-ARM-6X": 4,
    "NWR-LDR-360": 12,
    "NWR-DRV-BLDC": 61,
}

ROBOT_TELEMETRY: dict[str, dict[str, float | int]] = {
    "R-101": {"motor_temp_c": 41.2, "cycles": 18233},
    "R-102": {"motor_temp_c": 38.7, "cycles": 12004},
    "R-103": {"motor_temp_c": 44.9, "cycles": 25077},
}

DEMO_API_KEY = "dev-key-northwind"
RATE_LIMIT = 30
API_VERSION = "2026-09-01"


# --- Errors: RFC 9457 problem details (Ch. 7) ------------------------------

def problem(status: int, title: str, detail: str, type_: str = "about:blank") -> JSONResponse:
    """Build an RFC 9457 ``application/problem+json`` response."""
    return JSONResponse(
        status_code=status,
        media_type="application/problem+json",
        content={"type": type_, "title": title, "status": status, "detail": detail},
    )


# --- Platform state: rate limiting (Ch. 14) and idempotency (Ch. 10) -------

@dataclass
class PlatformState:
    """Mutable application state guarded by one lock (demo-scale)."""

    lock: threading.Lock = field(default_factory=threading.Lock)
    remaining: int = RATE_LIMIT
    orders: dict[str, dict[str, object]] = field(default_factory=dict)
    next_order_id: int = 1


class CursorError(ValueError):
    """Raised when a keyset pagination cursor cannot be decoded."""


# --- Schemas ----------------------------------------------------------------

class ProductOut(BaseModel):
    id: int
    sku: str
    name: str
    category: str
    unit_price_cents: int


class ProductPage(BaseModel):
    items: list[ProductOut]
    next_cursor: str | None = None
    api_version: str


class OrderIn(BaseModel):
    sku: str
    quantity: int = Field(ge=1, le=100)


class OrderOut(BaseModel):
    order_id: str
    sku: str
    quantity: int
    status: str
    api_version: str


def _encode_cursor(product_id: int) -> str:
    return base64.urlsafe_b64encode(f"cursor:{product_id}".encode()).decode()


def _decode_cursor(cursor: str) -> int:
    raw = base64.urlsafe_b64decode(cursor.encode()).decode()
    prefix, _, value = raw.partition(":")
    if prefix != "cursor" or not value.isdigit():
        raise ValueError(f"malformed cursor: {cursor!r}")
    return int(value)


# --- GraphQL read model (Ch. 12) -------------------------------------------

@strawberry.type
class CatalogItem:
    sku: str
    name: str
    category: str
    unit_price_cents: int
    availability: int


@strawberry.type
class RobotStatus:
    robot_id: str
    motor_temp_c: float
    cycles: int


@strawberry.type
class GraphQuery:
    @strawberry.field
    def catalog(self, category: str | None = None) -> list[CatalogItem]:
        rows = PRODUCTS
        if category is not None:
            rows = [p for p in rows if p["category"] == category]
        return [
            CatalogItem(
                sku=str(p["sku"]),
                name=str(p["name"]),
                category=str(p["category"]),
                unit_price_cents=int(p["unit_price_cents"]),
                availability=AVAILABILITY.get(str(p["sku"]), 0),
            )
            for p in rows
        ]

    @strawberry.field
    def robot_status(self, robot_id: str) -> RobotStatus | None:
        row = ROBOT_TELEMETRY.get(robot_id)
        if row is None:
            return None
        return RobotStatus(
            robot_id=robot_id,
            motor_temp_c=float(row["motor_temp_c"]),
            cycles=int(row["cycles"]),
        )


graphql_router = GraphQLRouter(strawberry.Schema(query=GraphQuery))


# --- Application factory -----------------------------------------------------

def create_app(state: PlatformState | None = None) -> FastAPI:
    """Build the platform app; inject state to keep tests isolated."""
    app = FastAPI(title="Northwind Robotics API Platform", version="1.0.0")
    platform = state or PlatformState()

    @app.middleware("http")
    async def platform_middleware(request: Request, call_next):  # type: ignore[no-untyped-def]
        if request.url.path in ("/healthz", "/docs", "/openapi.json"):
            return await call_next(request)
        if request.headers.get("X-API-Key") != DEMO_API_KEY:
            return problem(401, "Unauthorized", "missing or invalid X-API-Key header")
        with platform.lock:
            if platform.remaining <= 0:
                return JSONResponse(
                    status_code=429,
                    media_type="application/problem+json",
                    headers={"Retry-After": "30", "RateLimit": f'"limit={RATE_LIMIT}"'},
                    content={
                        "type": "about:blank",
                        "title": "Too Many Requests",
                        "status": 429,
                        "detail": "quota exhausted; retry after the Retry-After window",
                    },
                )
            platform.remaining -= 1
            remaining = platform.remaining
        response = await call_next(request)
        response.headers["RateLimit"] = f'"limit={RATE_LIMIT}, remaining={remaining}"'
        response.headers["RateLimit-Remaining"] = str(remaining)
        return response

    @app.exception_handler(RequestValidationError)
    async def validation_handler(
        _request: Request, exc: RequestValidationError
    ) -> JSONResponse:
        detail = "; ".join(
            f"{'.'.join(str(p) for p in err['loc'])}: {err['msg']}" for err in exc.errors()
        )
        return problem(422, "Unprocessable Entity", detail)

    # --- REST catalog (Ch. 6 + Ch. 8) -------------------------------------
    @app.get("/v1/products", response_model=ProductPage)
    def list_products(
        category: str | None = None,
        limit: int = Query(default=2, ge=1, le=50),
        cursor: str | None = None,
    ) -> ProductPage:
        rows = [p for p in PRODUCTS if category is None or p["category"] == category]
        rows.sort(key=lambda p: int(p["id"]))
        if cursor is not None:
            try:
                after = _decode_cursor(cursor)
            except ValueError as exc:
                raise CursorError(str(exc)) from exc
            rows = [p for p in rows if int(p["id"]) > after]
        page = rows[:limit]
        next_cursor = (
            _encode_cursor(int(page[-1]["id"])) if len(rows) > limit else None
        )
        return ProductPage(
            items=[ProductOut(**p) for p in page],  # type: ignore[arg-type]
            next_cursor=next_cursor,
            api_version=API_VERSION,
        )

    @app.exception_handler(CursorError)
    async def cursor_handler(_request: Request, exc: CursorError) -> JSONResponse:
        return problem(400, "Invalid Cursor", str(exc))

    # --- Orders with idempotency keys (Ch. 10) -----------------------------
    @app.post("/v1/orders", response_model=OrderOut, status_code=201)
    def create_order(order: OrderIn, request: Request) -> OrderOut | JSONResponse:
        key = request.headers.get("Idempotency-Key")
        if not key:
            return problem(
                400,
                "Missing Idempotency-Key",
                "POST /v1/orders requires an Idempotency-Key header",
            )
        with platform.lock:
            existing = platform.orders.get(key)
            if existing is not None:
                if existing["payload"] != order.model_dump():
                    return problem(
                        422,
                        "Idempotency Conflict",
                        "key reused with a different request payload",
                    )
                return OrderOut(
                    order_id=str(existing["order_id"]),
                    sku=order.sku,
                    quantity=order.quantity,
                    status="accepted",
                    api_version=API_VERSION,
                )
            if order.sku not in AVAILABILITY:
                return problem(404, "Unknown SKU", f"no catalog entry for {order.sku!r}")
            order_id = f"ORD-{platform.next_order_id:04d}"
            platform.next_order_id += 1
            platform.orders[key] = {
                "order_id": order_id,
                "payload": order.model_dump(),
            }
        return OrderOut(
            order_id=order_id,
            sku=order.sku,
            quantity=order.quantity,
            status="accepted",
            api_version=API_VERSION,
        )

    # --- SSE telemetry stream (Ch. 11) --------------------------------------
    @app.get("/v1/telemetry/stream")
    async def telemetry_stream() -> StreamingResponse:
        async def events() -> AsyncIterator[str]:
            event_id = 0
            for robot_id in sorted(ROBOT_TELEMETRY):
                row = ROBOT_TELEMETRY[robot_id]
                event_id += 1
                yield (
                    f"id: {event_id}\n"
                    f"event: robot-reading\n"
                    f'data: {{"robot_id": "{robot_id}", '
                    f'"motor_temp_c": {row["motor_temp_c"]}, '
                    f'"cycles": {row["cycles"]}}}\n\n'
                )

        return StreamingResponse(
            events(),
            media_type="text/event-stream",
            headers={"Cache-Control": "no-cache"},
        )

    @app.get("/healthz")
    def healthz() -> dict[str, str]:
        return {"status": "ok"}

    app.include_router(graphql_router, prefix="/graphql")
    return app


if __name__ == "__main__":
    uvicorn.run(create_app(), host="127.0.0.1", port=8121)
\end{minted}

The smoke suite proves the surfaces work \emph{together} --- through one
middleware stack, with one state --- using an in-process ASGI client, so no
sockets, servers, or network are involved:

\begin{minted}{python}
"""In-process smoke tests for the Northwind platform slice.

Every request runs through httpx + ASGITransport: no sockets, no
network, fully deterministic. The suite proves the four surfaces
(REST, GraphQL, SSE, problem-details errors) work together on one app.
"""

from __future__ import annotations

import json
from collections.abc import Iterator

import pytest
from fastapi.testclient import TestClient

from northwind_platform import DEMO_API_KEY, create_app

AUTH = {"X-API-Key": DEMO_API_KEY}


@pytest.fixture()
def client() -> Iterator[TestClient]:
    """Fresh app + in-process client per test (isolated rate limiter)."""
    with TestClient(create_app()) as c:
        yield c


# --- Auth and rate limiting --------------------------------------------------

def test_missing_api_key_is_problem_json(client: TestClient) -> None:
    r = client.get("/v1/products")
    assert r.status_code == 401
    assert r.headers["content-type"] == "application/problem+json"
    body = r.json()
    assert body["status"] == 401
    assert {"type", "title", "status", "detail"} <= body.keys()


def test_rate_limit_headers_present(client: TestClient) -> None:
    r = client.get("/v1/products", headers=AUTH)
    assert r.status_code == 200
    assert "RateLimit" in r.headers
    assert int(r.headers["RateLimit-Remaining"]) < 30


def test_rate_limit_exhaustion_returns_429() -> None:
    with TestClient(create_app()) as c:
        last = None
        for _ in range(31):
            last = c.get("/v1/products", headers=AUTH)
        assert last is not None and last.status_code == 429
        assert last.headers["Retry-After"] == "30"


# --- REST catalog ------------------------------------------------------------

def test_product_page_one_and_cursor(client: TestClient) -> None:
    page1 = client.get("/v1/products", headers=AUTH, params={"limit": 2}).json()
    assert len(page1["items"]) == 2
    assert page1["next_cursor"] is not None
    assert page1["api_version"]

    page2 = client.get(
        "/v1/products",
        headers=AUTH,
        params={"limit": 2, "cursor": page1["next_cursor"]},
    ).json()
    assert len(page2["items"]) == 1
    assert page2["next_cursor"] is None
    ids = [p["id"] for p in page1["items"] + page2["items"]]
    assert ids == [1, 2, 3]


def test_bad_cursor_is_problem_400(client: TestClient) -> None:
    r = client.get(
        "/v1/products", headers=AUTH, params={"cursor": "not-a-cursor"}
    )
    assert r.status_code == 400
    assert r.json()["title"] == "Invalid Cursor"


def test_validation_error_is_problem_422(client: TestClient) -> None:
    r = client.get("/v1/products", headers=AUTH, params={"limit": 0})
    assert r.status_code == 422
    assert r.headers["content-type"] == "application/problem+json"


# --- Orders + idempotency -----------------------------------------------------

def test_order_requires_idempotency_key(client: TestClient) -> None:
    r = client.post("/v1/orders", headers=AUTH, json={"sku": "NWR-ARM-6X", "quantity": 1})
    assert r.status_code == 400


def test_order_replay_returns_same_order_id(client: TestClient) -> None:
    headers = {**AUTH, "Idempotency-Key": "smoke-key-1"}
    payload = {"sku": "NWR-LDR-360", "quantity": 2}
    first = client.post("/v1/orders", headers=headers, json=payload)
    second = client.post("/v1/orders", headers=headers, json=payload)
    assert first.status_code == 201
    assert second.json()["order_id"] == first.json()["order_id"]


def test_idempotency_key_conflict_is_422(client: TestClient) -> None:
    headers = {**AUTH, "Idempotency-Key": "smoke-key-2"}
    client.post("/v1/orders", headers=headers, json={"sku": "NWR-ARM-6X", "quantity": 1})
    conflict = client.post(
        "/v1/orders", headers=headers, json={"sku": "NWR-ARM-6X", "quantity": 9}
    )
    assert conflict.status_code == 422
    assert conflict.json()["title"] == "Idempotency Conflict"


def test_unknown_sku_is_problem_404(client: TestClient) -> None:
    headers = {**AUTH, "Idempotency-Key": "smoke-key-3"}
    r = client.post("/v1/orders", headers=headers, json={"sku": "NOPE-1", "quantity": 1})
    assert r.status_code == 404
    assert r.json()["status"] == 404


# --- GraphQL read model --------------------------------------------------------

def test_graphql_catalog_joins_availability(client: TestClient) -> None:
    query = "{ catalog { sku name availability } }"
    r = client.post("/graphql", headers=AUTH, json={"query": query})
    assert r.status_code == 200
    rows = {item["sku"]: item for item in r.json()["data"]["catalog"]}
    assert rows["NWR-LDR-360"]["availability"] == 12


def test_graphql_robot_status(client: TestClient) -> None:
    query = '{ robotStatus(robotId: "R-101") { robotId motorTempC cycles } }'
    r = client.post("/graphql", headers=AUTH, json={"query": query})
    assert r.status_code == 200
    status = r.json()["data"]["robotStatus"]
    assert status["robotId"] == "R-101"
    assert status["motorTempC"] == pytest.approx(41.2)


# --- SSE telemetry --------------------------------------------------------------

def test_sse_stream_yields_three_robot_events(client: TestClient) -> None:
    with client.stream("GET", "/v1/telemetry/stream", headers=AUTH) as r:
        assert r.status_code == 200
        assert r.headers["content-type"].startswith("text/event-stream")
        body = "".join(r.iter_text())
    events = [chunk for chunk in body.strip().split("\n\n") if chunk]
    assert len(events) == 3
    assert all(chunk.startswith("id: ") for chunk in events)
    payload = events[0].split("data: ", 1)[1]
    assert json.loads(payload)["robot_id"] == "R-101"
\end{minted}

\subsection{Listing 3 --- The Architecture Conformance Checker}
\label{sec:code-checker}

Chapter~\ref{ch:lifecycle} argued that governance must be executable or it
is decorative. This checker is the proof: it loads every OpenAPI
\cite{openapi31} and AsyncAPI \cite{asyncapi} manifest in a folder and
verifies the platform's invariants --- security schemes with descriptions
and a global \texttt{security} requirement (O-001), RFC~9457 error shapes on
every 4xx/5xx response (O-002) \cite{rfc9457}, keyset pagination conventions
on collection GETs (O-003), date-based version identifiers (O-004), and the
event envelope on AsyncAPI publish payloads (A-001). Local
\texttt{\$ref} resolution keeps the checker honest against real spec style.

\begin{minted}{python}
"""Architecture conformance checker (Chapter 21).

Turns the book's governance rules into executable policy. Given a
manifests folder of OpenAPI (YAML/JSON) and AsyncAPI documents, the
checker verifies the platform invariants that Chapter 20 establishes
as linted, CI-enforced rules:

* O-001: every spec declares a security scheme with a non-empty
  ``description`` and applies ``security`` globally.
* O-002: every 4xx/5xx response references the shared
  ``ProblemDetails`` component (RFC 9457).
* O-003: every collection GET documents ``limit`` and ``cursor``
  query parameters and a ``nextCursor`` response field.
* O-004: the ``info.version`` header version matches
  ``YYYY-MM-DD`` (the platform's date-based versioning rule).
* A-001: every AsyncAPI publish payload message defines ``event_id``
  and ``occurred_at`` (the platform's event envelope).

Offline, deterministic. Requires only PyYAML. Python 3.11+.
"""

from __future__ import annotations

import re
import sys
from dataclasses import dataclass
from pathlib import Path
from typing import Any

import yaml

DATE_VERSION = re.compile(r"^\d{4}-\d{2}-\d{2}$")
ENVELOPE_FIELDS = ("event_id", "occurred_at")


@dataclass(frozen=True)
class Violation:
    """A single governance rule breach in one manifest."""

    rule: str
    spec: str
    detail: str

    def render(self) -> str:
        return f"[{self.rule}] {self.spec}: {self.detail}"


def _load_specs(folder: Path) -> dict[str, Any]:
    specs: dict[str, Any] = {}
    for path in sorted(folder.rglob("*")):
        if path.suffix.lower() in {".yaml", ".yml", ".json"}:
            specs[str(path.relative_to(folder))] = yaml.safe_load(
                path.read_text(encoding="utf-8")
            )
    return specs


def _resolve_local_ref(document: Any, node: Any, depth: int = 0) -> Any:
    """Resolve one local ``#/...`` reference against its own document."""
    if depth > 10 or not isinstance(node, dict) or "$ref" not in node:
        return node
    ref = str(node["$ref"])
    if not ref.startswith("#/"):
        return node  # external refs are out of scope for the offline checker
    target: Any = document
    for part in ref[2:].split("/"):
        if not isinstance(target, dict) or part not in target:
            return node
        target = target[part]
    return _resolve_local_ref(document, target, depth + 1)


def _schema_has_property(document: Any, schema: Any, name: str) -> bool:
    """True when ``name`` appears in the schema or any composed/ref'd part."""
    schema = _resolve_local_ref(document, schema)
    if not isinstance(schema, dict):
        return False
    if name in (schema.get("properties") or {}):
        return True
    for key in ("allOf", "anyOf", "oneOf"):
        if any(
            _schema_has_property(document, part, name) for part in schema.get(key, [])
        ):
            return True
    return False


def _collect_refs(document: Any, schema: Any, depth: int = 0) -> list[str]:
    """Collect every ``$ref`` target reachable from a schema (any depth)."""
    schema = _resolve_local_ref(document, schema)
    if depth > 10 or not isinstance(schema, dict):
        return []
    refs: list[str] = []
    raw = schema.get("$ref")
    if isinstance(raw, str):
        refs.append(raw)
    for key in ("allOf", "anyOf", "oneOf"):
        for part in schema.get(key, []):
            refs.extend(_collect_refs(document, part, depth + 1))
    return refs


def _content_schemas(content: Any) -> list[Any]:
    if not isinstance(content, dict):
        return []
    return [media.get("schema") for media in content.values() if isinstance(media, dict)]


def _check_openapi(name: str, spec: Any) -> list[Violation]:
    violations: list[Violation] = []
    if not isinstance(spec, dict) or "openapi" not in spec:
        return violations

    # O-001: auth scheme presence and global security requirement.
    schemes = (
        (spec.get("components") or {}).get("securitySchemes") or {}
    )
    if not schemes:
        violations.append(Violation("O-001", name, "no securitySchemes defined"))
    for scheme_name, scheme in schemes.items():
        if not (scheme or {}).get("description"):
            violations.append(
                Violation("O-001", name, f"scheme {scheme_name!r} lacks a description")
            )
    if not spec.get("security"):
        violations.append(
            Violation("O-001", name, "no global security requirement applied")
        )

    # O-002: error responses must reference ProblemDetails.
    paths = spec.get("paths") or {}
    for path, item in paths.items():
        for method, operation in (item or {}).items():
            if method not in ("get", "post", "put", "patch", "delete"):
                continue
            responses = (operation or {}).get("responses") or {}
            for status, response in responses.items():
                if not str(status).startswith(("4", "5")):
                    continue
                schemas = _content_schemas((response or {}).get("content"))
                refs: list[str] = []
                for schema in schemas:
                    raw = (schema or {}).get("$ref") if isinstance(schema, dict) else None
                    if raw:
                        refs.append(raw)
                    else:
                        refs.extend(_collect_refs(spec, schema))
                if not any(r.endswith("/ProblemDetails") for r in refs):
                    violations.append(
                        Violation(
                            "O-002",
                            name,
                            f"{method.upper()} {path} {status} does not reference "
                            "ProblemDetails",
                        )
                    )

    # O-003: collection GETs must document limit/cursor pagination.
    for path, item in paths.items():
        get_op = (item or {}).get("get")
        if not get_op or not path.rstrip("/").endswith("s"):
            continue
        params = {
            (p or {}).get("name")
            for p in get_op.get("parameters") or []
            if isinstance(p, dict) and p.get("in") == "query"
        }
        missing = {"limit", "cursor"} - params
        if missing:
            violations.append(
                Violation(
                    "O-003",
                    name,
                    f"GET {path} missing query params: {sorted(missing)}",
                )
            )
        ok_schemas = _content_schemas(
            ((get_op.get("responses") or {}).get("200") or {}).get("content")
        )
        if not any(
            _schema_has_property(spec, s, "nextCursor") for s in ok_schemas
        ):
            violations.append(
                Violation("O-003", name, f"GET {path} 200 schema lacks nextCursor")
            )

    # O-004: date-based version header convention.
    version = str((spec.get("info") or {}).get("version", ""))
    if not DATE_VERSION.match(version):
        violations.append(
            Violation(
                "O-004",
                name,
                f"info.version {version!r} is not a YYYY-MM-DD date version",
            )
        )
    return violations


def _check_asyncapi(name: str, spec: Any) -> list[Violation]:
    violations: list[Violation] = []
    if not isinstance(spec, dict) or "asyncapi" not in spec:
        return violations
    channels = spec.get("channels") or {}
    for channel, item in channels.items():
        publish = (item or {}).get("publish")
        if not publish:
            continue
        message = _resolve_local_ref(spec, publish.get("message") or {})
        payload = _resolve_local_ref(spec, (message or {}).get("payload") or {})
        missing = [
            f for f in ENVELOPE_FIELDS if not _schema_has_property(spec, payload, f)
        ]
        if missing:
            violations.append(
                Violation(
                    "A-001",
                    name,
                    f"channel {channel!r} publish payload missing envelope "
                    f"fields: {missing}",
                )
            )
    return violations


def check_folder(folder: Path) -> list[Violation]:
    """Check every spec under ``folder`` and return all violations."""
    specs = _load_specs(folder)
    if not specs:
        raise FileNotFoundError(f"no YAML/JSON manifests found under {folder}")
    violations: list[Violation] = []
    for name, spec in specs.items():
        violations.extend(_check_openapi(name, spec))
        violations.extend(_check_asyncapi(name, spec))
    return violations


def main(argv: list[str]) -> int:
    folder = Path(argv[1]) if len(argv) > 1 else Path("manifests")
    try:
        violations = check_folder(folder)
    except FileNotFoundError as exc:
        print(f"error: {exc}", file=sys.stderr)
        return 2
    if violations:
        print(f"FAIL: {len(violations)} governance violation(s)")
        for v in violations:
            print(f"  {v.render()}")
        return 1
    print(f"PASS: all governance invariants hold under {folder}")
    return 0


if __name__ == "__main__":
    raise SystemExit(main(sys.argv))
\end{minted}

The checker is validated against two synthetic manifest folders shipped with
the lab: a compliant Northwind folder and a legacy folder that breaches every
rule family. A dedicated test proves that \texttt{\$ref}-based AsyncAPI
payloads resolve instead of false-positiving.

\begin{minted}{python}
"""Tests for the architecture conformance checker.

Runs against two synthetic fixture folders: ``fixtures_pass`` (a
compliant Northwind manifests folder) and ``fixtures_fail`` (a legacy
folder that breaches every rule family).
"""

from __future__ import annotations

from pathlib import Path

import pytest

from conformance_checker import check_folder, main

HERE = Path(__file__).parent
PASS_DIR = HERE / "fixtures_pass"
FAIL_DIR = HERE / "fixtures_fail"


def test_compliant_folder_passes() -> None:
    assert check_folder(PASS_DIR) == []


def test_legacy_folder_fails() -> None:
    violations = check_folder(FAIL_DIR)
    assert violations, "legacy folder must produce violations"


def test_every_rule_family_fires_on_legacy_folder() -> None:
    rules = {v.rule for v in check_folder(FAIL_DIR)}
    assert {"O-001", "O-002", "O-003", "O-004", "A-001"} <= rules


def test_main_exit_codes() -> None:
    assert main(["prog", str(PASS_DIR)]) == 0
    assert main(["prog", str(FAIL_DIR)]) == 1


def test_missing_folder_exit_code() -> None:
    assert main(["prog", str(HERE / "no_such_dir")]) == 2


def test_checker_is_deterministic() -> None:
    assert check_folder(FAIL_DIR) == check_folder(FAIL_DIR)


def test_asyncapi_ref_payload_is_resolved(tmp_path: Path) -> None:
    """A publish payload referenced via components must not false-positive."""
    spec = tmp_path / "ref.asyncapi.yaml"
    spec.write_text(
        "\n".join(
            [
                'asyncapi: "2.6.0"',
                "info: { title: Ref Events, version: '1.0.0' }",
                "channels:",
                "  robot.readings:",
                "    publish:",
                "      message: { $ref: '#/components/messages/Reading' }",
                "components:",
                "  messages:",
                "    Reading:",
                "      payload: { $ref: '#/components/schemas/Envelope' }",
                "  schemas:",
                "    Envelope:",
                "      type: object",
                "      properties:",
                "        event_id: { type: string }",
                "        occurred_at: { type: string, format: date-time }",
            ]
        ),
        encoding="utf-8",
    )
    assert check_folder(tmp_path) == []


if __name__ == "__main__":
    raise SystemExit(pytest.main([__file__, "-q"]))
\end{minted}

The compliant manifests folder (abridged here to the catalog spec; the
AsyncAPI telemetry spec follows the same shape):

\begin{minted}{yaml}
openapi: 3.1.0
info:
  title: Northwind Catalog API
  version: "2026-09-01"
security:
  - ApiKeyAuth: []
paths:
  /v1/products:
    get:
      summary: List catalog products
      parameters:
        - name: limit
          in: query
          schema: { type: integer, minimum: 1, maximum: 50, default: 20 }
        - name: cursor
          in: query
          schema: { type: string }
      responses:
        "200":
          description: Page of products
          content:
            application/json:
              schema:
                type: object
                properties:
                  items:
                    type: array
                    items: { $ref: "#/components/schemas/Product" }
                  nextCursor: { type: string }
        "400":
          description: Invalid cursor
          content:
            application/problem+json:
              schema: { $ref: "#/components/schemas/ProblemDetails" }
        "401":
          description: Missing or invalid API key
          content:
            application/problem+json:
              schema: { $ref: "#/components/schemas/ProblemDetails" }
components:
  securitySchemes:
    ApiKeyAuth:
      type: apiKey
      in: header
      name: X-API-Key
      description: Partner API key issued via the Northwind developer portal.
  schemas:
    Product:
      type: object
      properties:
        id: { type: integer }
        sku: { type: string }
        name: { type: string }
        category: { type: string }
        unit_price_cents: { type: integer }
    ProblemDetails:
      type: object
      properties:
        type: { type: string }
        title: { type: string }
        status: { type: integer }
        detail: { type: string }
\end{minted}

\section{Common Pitfalls \& Anti-Patterns}

The failure modes below are synthesis-level: they appear only when you
combine everything, which is exactly why they are common at platform scale.

\subsection{Resume-Driven Protocol Selection}

Choosing GraphQL or gRPC because it is fashionable, then discovering the
public partner audience cannot consume it, is the canonical platform mistake.
The decision card exists precisely to make this failure awkward: if you
cannot fill in the six criteria, you have not done the analysis. The repair
--- a BFF shim, a gRPC-Web proxy, a parallel REST surface --- always costs
more than the analysis would have.

\subsection{One-Protocol-Fits-All Mandates}

The mirror image of fashion is the organizational mandate: ``everything is
REST'' or ``everything is gRPC.'' Mandates feel like governance but are
actually the refusal to govern. The reference architecture of
Section~\ref{sec:ref-arch} runs four profiles simultaneously \emph{because
the constraints differ per consumer}; a mandate forces the worst profile onto
at least one consumer class. Govern invariants (errors, auth, pagination,
observability), not protocols.

\subsection{Big-Bang Platform Builds}

Attempting all four milestones of the Grand Capstone in one heroic push
produces a platform that is 90\,\% built and 0\,\% operable. The phased plan
of Table~\ref{tab:milestones} is ordered so that each milestone is
independently shippable and teaches the organization something the next one
needs. If M1's REST core cannot pass its contract tests, M3's event pipeline
will not save you.

\subsection{Skipping Governance Until Fifty APIs Exist}

Governance feels premature with three APIs and impossible with fifty. The
cheap moment is the beginning: three rules (error shape, auth scheme,
pagination convention) in a conformance checker like
Listing~3 (Section~\ref{sec:code-checker}) cost an afternoon when you have three specs and a
quarter when you have fifty --- and retrofitting means negotiating with fifty
teams' shipped contracts \cite{gehl2021apidesign}.

\subsection{Adopting Frontiers Without Operational Maturity}

HTTP/3, edge runtimes, and MCP tool surfaces amplify whatever you already
are. Teams without SLOs, traces, and rollback discipline do not get faster
from QUIC; they get \emph{less debuggable}. Table~\ref{tab:frontiers} lists a
prerequisite per frontier; treat it as a gate, not a suggestion.

\subsection{Synthesis Without Ownership}

A reference architecture diagram that no team owns decays into a poster.
Every box in Figure~\ref{fig:ref-arch} must have a named owning team, an SLO,
and a runbook, or the architecture is fiction
\cite{newman2021microservices}.

\section{Hands-On Production Lab \& Real-World Case Study}

\subsection{Lab: Run the Full Toolkit}

\textbf{Goal:} exercise all three listings end to end, offline, in under ten
minutes.

\textbf{Steps (PowerShell):}

\begin{minted}{powershell}
cd "code\ch21"
.\.venv\Scripts\Activate.ps1

# 1. Prove all 30 tests pass (advisor + platform + checker).
python -m pytest -q

# 2. Run the decision advisor against the five canonical scenarios.
python protocol_advisor.py

# 3. Check governance: compliant folder passes, legacy folder fails.
python conformance_checker.py fixtures_pass
python conformance_checker.py fixtures_fail

# 4. Boot the platform slice and probe every surface.
python northwind_platform.py   # serves http://127.0.0.1:8121
\end{minted}

In a second PowerShell window:

\begin{minted}{powershell}
$h = @{ "X-API-Key" = "dev-key-northwind" }

# REST catalog with keyset pagination (note RateLimit-* headers).
curl.exe -s -D - "http://127.0.0.1:8121/v1/products?limit=2" -H "X-API-Key: dev-key-northwind"

# Problem-details error: no auth header.
curl.exe -s "http://127.0.0.1:8121/v1/products"

# Idempotent order placement: run twice, compare order_id.
$body = '{"sku":"NWR-LDR-360","quantity":2}'
curl.exe -s -X POST "http://127.0.0.1:8121/v1/orders" -H "X-API-Key: dev-key-northwind" -H "Idempotency-Key: lab-1" -H "Content-Type: application/json" -d $body
curl.exe -s -X POST "http://127.0.0.1:8121/v1/orders" -H "X-API-Key: dev-key-northwind" -H "Idempotency-Key: lab-1" -H "Content-Type: application/json" -d $body

# GraphQL read model.
$q = '{"query":"{ catalog { sku name availability } }"}'
curl.exe -s -X POST "http://127.0.0.1:8121/graphql" -H "X-API-Key: dev-key-northwind" -H "Content-Type: application/json" -d $q

# SSE telemetry stream (three robot events, then the demo stream ends).
curl.exe -s -N "http://127.0.0.1:8121/v1/telemetry/stream" -H "X-API-Key: dev-key-northwind"
\end{minted}

\textbf{Expected results:} 30 tests pass; the advisor prints REST, gRPC,
GraphQL, event-driven, and event-driven for scenarios A--E with rationale
traces; the checker prints \texttt{PASS} for \texttt{fixtures\_pass} and
seven violations for \texttt{fixtures\_fail}; the live server returns a
paginated product page with \texttt{RateLimit-Remaining}, an
\texttt{application/problem+json} 401 without the key, the same
\texttt{order\_id} on both POSTs, three catalog rows from GraphQL with
joined availability, and three \texttt{robot-reading} SSE events.

\textbf{Extensions:} (a) flip one scenario's constraints in the advisor
(e.g., make scenario~C cacheable) and watch the recommendation change;
(b) add a violating spec to \texttt{fixtures\_fail} and predict the rule
codes before running; (c) exhaust the 30-request rate limit against the live
server and confirm the 429 carries \texttt{Retry-After}.

\subsection{Case Study: Three Years of an Enterprise API Platform}

\textbf{Year 1 --- the product year.} Northwind Robotics ships a public REST
catalog and orders API to resellers: OpenAPI-first, RFC~9457 errors from day
one, keyset pagination, idempotency keys on orders. The team debates GraphQL
for the partner audience and rejects it with the decision card: unknown third
parties plus a cacheable catalog is a REST scenario, full stop. A two-rule
conformance linter (error shape, auth scheme) runs in CI from the first
spec. Year-end incident count: one --- a pagination drift bug caught by a
contract test, not a partner.

\textbf{Year 2 --- the platform year.} The mobile app launches, and its
nested screens justify a GraphQL read model (scenario~C) with depth limits
and persisted queries. Inventory moves to internal gRPC when the checkout hot
path misses its latency budget under REST-over-JSON. Order events go through
a transactional outbox to signed partner webhooks after one painful week in
which a dual-write bug (database committed, webhook never sent) teaches the
team why Chapter~\ref{ch:async_apis} insists on the outbox
\cite{kleppmann2017ddia}. Rate tiers launch with IETF headers; the noisy
top-1\,\% integrators get quota offers instead of silent 429s. The
conformance checker, now five rules, blocks two specs that would have shipped
undocumented breaking changes.

\textbf{Year 3 --- the strategy year.} OTel traces span edge to outbox; SLOs
with error budgets replace uptime folklore \cite{beyer2016sre}. The platform
team evaluates HTTP/3 (adopted free at the CDN edge; internal gRPC stays on
HTTP/2), runs a pilot exposing read-only catalog tools via MCP to an internal
support copilot with tool schemas generated from the same OpenAPI source, and
declines an edge-compute rewrite of the order write path --- regional
affinity could not satisfy the single-region consistency requirement, a
\texttt{wait} decision the frontier table predicts. The year's most quoted
internal document is not a design doc but the conformance checker's violation
log: governance-as-code turned fifty APIs into one coherent platform.

\textbf{Lessons:} protocols were chosen per consumer and revisited only when
constraints changed; governance was cheapest on day one; the outbox and
idempotency investments paid out every incident; and every frontier decision
had an explicit adopt/trial/wait rationale tied to measured constraints.

\section{Self-Assessment Exercises \& Deep-Dive Questions}

Work these without rereading the chapter; then check your reasoning against
the decision card and the reference architecture.

\subsection{Concept Review}

\begin{enumerate}
  \item Why is the delivery model the \emph{first} question in the decision
        tree rather than the consumer type? Give an example where ordering
        them the other way produces a wrong recommendation.
  \item The comparison matrix calls its entries ``tendencies, not laws.''
        Pick one cell and construct a legitimate exception.
  \item In the reference architecture, why does the outbox relay read from
        the database instead of the service publishing to the broker
        directly? What exactly can go wrong otherwise
        \cite{kleppmann2017ddia}?
  \item Which single header pair makes a public REST catalog dramatically
        cheaper to operate, and why do GraphQL endpoints not benefit equally
        \cite{rfc9110}?
  \item The conformance checker's O-002 rule requires a
        \texttt{ProblemDetails} reference on every 4xx/5xx response. What
        failure class of Table~\ref{tab:failure-taxonomy} does this prevent,
        and for whom?
  \item Why does the advisor break ties deterministically instead of
        reporting ``either is fine''? What property of the test suite
        depends on it?
  \item Explain why MCP tool schemas should be \emph{generated from} the
        OpenAPI source of truth rather than maintained by hand.
  \item The case study's year-3 team declined an edge-compute rewrite of the
        order path. Reconstruct their reasoning using
        Table~\ref{tab:frontiers}.
\end{enumerate}

\subsection{Architecture-Design Drills}

Each drill gives a consumer, a constraint set, and asks for a defended
design. Use the decision card, name the profile(s), sketch the surfaces, and
state the top two failure modes with mitigations.

\begin{enumerate}
  \item \textbf{Drill 1 --- Hospital device monitoring.} Bedside devices
        emit vitals every second; nurses view live dashboards; auditors pull
        historical ranges. Constraints: no public internet exposure, p99
        dashboard freshness $<$ 2\,s, audit reads cacheable for 24\,h,
        small disciplined team. Design all three surfaces.
  \item \textbf{Drill 2 --- Marketplace payout API.} Sellers trigger
        payouts; duplicates are catastrophic; partner banks are slow and
        flaky. Constraints: public consumers, strict exactly-once semantics
        at the boundary, audit trail, 100 requests/min typical with 10$\times$
        month-end bursts. Which chapters' patterns combine here?
  \item \textbf{Drill 3 --- Multi-tenant analytics backend.} A SPA renders
        arbitrary user-configured dashboards over 40 entity types.
        Constraints: payload shape unknowable at design time, metered mobile
        variant planned next year, team of four. Defend or reject GraphQL
        using the card, and design the guardrails whichever way you land.
  \item \textbf{Drill 4 --- Fleet firmware update orchestration.} 50\,000
        devices must receive signed update commands and report progress;
        connectivity is intermittent. Constraints: server-initiated delivery,
        resume across disconnects, no inbound connections to devices allowed
        on some networks. Choose between SSE, WebSockets, gRPC streaming,
        and polling-with-ETags --- and defend the choice.
  \item \textbf{Drill 5 --- Your own platform.} Take a system you know
        professionally. Fill in the six criteria for each of its API
        surfaces. Does the card reproduce the choices that were actually
        made? Where it does not, was the card wrong or was the choice?
\end{enumerate}

\section{Interview Q\&A Vault}

Thirty-two synthesis-level questions, ordered junior $\rightarrow$ staff.
Answers are deliberately terse; the chapters cited hold the depth.

\subsection*{Junior}

\begin{enumerate}[label=\textbf{Q\arabic*.}, leftmargin=3em]
  \item \textbf{When would you choose REST over GraphQL for a new public
        API?} --- When consumers are unknown third parties, payloads are
        fixed resources, and CDN caching matters: REST's addressable GETs are
        the only profile that exploits shared HTTP caches \cite{rfc9110}, and
        OpenAPI tooling is universally understood \cite{openapi31}.
        (Ch.~\ref{ch:rest_style})
  \item \textbf{What is RFC~9457 and why should every error use it?} ---
        A standard JSON problem-details shape (\texttt{type},
        \texttt{title}, \texttt{status}, \texttt{detail}) so clients handle
        errors programmatically instead of parsing prose \cite{rfc9457}.
        (Ch.~\ref{ch:problem_details})
  \item \textbf{Offset or keyset pagination --- and why?} --- Keyset:
        stable under concurrent inserts, $O(1)$ deep pages, and it composes
        with cursors; offset drifts and scans \cite{rfc8288}.
        (Ch.~\ref{ch:pagination})
  \item \textbf{What does an idempotency key buy you?} --- Safe retries:
        the server deduplicates replays of the same key, so a flaky network
        cannot double-charge or double-order. (Ch.~\ref{ch:idempotency})
  \item \textbf{REST vs.\ gRPC in one sentence?} --- REST optimizes for
        reach and cacheability over standard HTTP; gRPC optimizes for typed,
        low-latency service-to-service RPC over HTTP/2
        \cite{rfc9113,grpcdocs}. (Ch.~\ref{ch:grpc})
  \item \textbf{What headers belong on a rate-limited response?} ---
        \texttt{429} with \texttt{Retry-After} plus the IETF
        \texttt{RateLimit} family so well-behaved clients can self-throttle
        \cite{ietf_ratelimit_headers}. (Ch.~\ref{ch:rate_limiting})
  \item \textbf{Why pin timeouts on every outbound call?} --- An unbounded
        wait turns a slow dependency into your outage; timeouts plus retries
        with backoff bound the blast radius \cite{nygard2018releaseit}.
        (Ch.~\ref{ch:consuming_python})
  \item \textbf{What is the OpenAPI document's role in a team?} --- The
        machine-readable source of truth for the contract: docs, tests,
        lint, and SDKs all derive from it \cite{openapi31}.
        (Ch.~\ref{ch:lifecycle})
\end{enumerate}

\subsection*{Mid}

\begin{enumerate}[label=\textbf{Q\arabic*.}, leftmargin=3em, start=9]
  \item \textbf{How do you version an API without breaking consumers?} ---
        Additive changes within a version; date-based versions for breaking
        ones; \texttt{Sunset} headers and a published deprecation policy
        \cite{rfc8594}. (Ch.~\ref{ch:versioning})
  \item \textbf{Explain the transactional outbox in ninety seconds.} ---
        Write the state change \emph{and} the event row in one transaction;
        a relay publishes from the outbox; consumers get at-least-once
        delivery and the dual-write anomaly disappears
        \cite{kleppmann2017ddia}. (Ch.~\ref{ch:async_apis})
  \item \textbf{When is GraphQL the wrong answer?} --- Cacheable public
        reads, teams without schema-governance maturity, and payloads that
        are actually fixed resources; you pay resolver complexity and lose
        HTTP caching for nothing \cite{graphqlspec}.
        (Ch.~\ref{ch:graphql})
  \item \textbf{How do you secure a webhook?} --- HMAC signatures with
        rotated secrets, timestamp tolerance against replay, delivery logs,
        and a replay endpoint; treat inbound signature failures as security
        events \cite{owaspapisecurity}. (Ch.~\ref{ch:async_apis})
  \item \textbf{SSE or WebSockets?} --- SSE when the flow is server-to-client
        text events: it rides plain HTTP, auto-reconnects, and resumes via
        \texttt{Last-Event-ID}; WebSockets when you need bidirectional or
        binary. (Ch.~\ref{ch:async_apis})
  \item \textbf{What belongs in a health endpoint vs.\ a readiness
        endpoint?} --- Liveness: process is up. Readiness: dependencies
        (database, broker) verified, safe to route traffic. Conflating them
        turns dependency blips into restart storms.
        (Ch.~\ref{ch:deployment})
  \item \textbf{How do circuit breakers and retries interact?} --- Retries
        handle transient faults; the breaker stops hammering a down
        dependency; retry budgets prevent retry storms from becoming a
        self-inflicted DDoS \cite{nygard2018releaseit}.
        (Ch.~\ref{ch:microservices})
  \item \textbf{Walk through what the gateway should and should not do.} ---
        Should: auth offload \cite{rfc6749,rfc7519}, rate tiers, correlation
        IDs, request validation. Should not: business logic, data
        transformation that services then cannot evolve independently.
        (Ch.~\ref{ch:microservices}, Ch.~\ref{ch:rate_limiting})
\end{enumerate}

\subsection*{Senior}

\begin{enumerate}[label=\textbf{Q\arabic*.}, leftmargin=3em, start=17]
  \item \textbf{Design the API surface for a real-time bidding system under
        a 50\,ms budget.} --- gRPC streaming between participants and the
        auction core for latency; events for post-trade fan-out; REST only
        for management; the decision card lands here via consumer=internal,
        latency=extreme. (Ch.~\ref{ch:grpc}, this chapter)
  \item \textbf{When would you expose gRPC publicly?} --- Rarely, and only
        for machine consumers you influence: SDK-shipped partners, not
        browsers. Budget for a gRPC-Web/JSON transcoding proxy, dual
        documentation, and firewall pain; if any of those is unacceptable,
        it was a REST API. (Ch.~\ref{ch:grpc})
  \item \textbf{How do you make four protocol stacks feel like one
        platform?} --- Govern the invariants, not the protocols: one error
        model (RFC~9457 \cite{rfc9457}), one auth story, one pagination
        convention, one tracing context, one portal --- enforced by
        executable policy. (Ch.~\ref{ch:lifecycle}, this chapter)
  \item \textbf{Your GraphQL gateway is melting under expensive queries.
        Now what?} --- Persisted queries, depth and cost limits, dataloaders
        for N+1, per-field tracing to find the hot resolvers, and caching at
        the edge only for truly public anonymous queries \cite{graphqlspec}.
        (Ch.~\ref{ch:graphql}, Ch.~\ref{ch:observability})
  \item \textbf{How do you roll out a breaking change to 200 partners?} ---
        You do not: you add the new shape alongside, emit \texttt{Sunset}
        and \texttt{Link} headers \cite{rfc8594,rfc8288}, measure old-version
        traffic, contact the long tail, and only then schedule removal.
        (Ch.~\ref{ch:versioning})
  \item \textbf{What does your error budget buy you operationally?} --- A
        neutral throttle: budget remaining $\Rightarrow$ ship; budget burned
        $\Rightarrow$ freeze features and invest in reliability
        \cite{beyer2016sre}. It converts reliability arguments into
        arithmetic. (Ch.~\ref{ch:observability})
  \item \textbf{Design the telemetry pipeline for 3\,000 robots.} ---
        Event-driven ingestion to a broker with envelope fields
        (\texttt{event\_id}, \texttt{occurred\_at}), consumer-side
        idempotency, SSE fan-out to dashboards, and AsyncAPI contracts
        linted in CI \cite{asyncapi}. (Ch.~\ref{ch:async_apis})
  \item \textbf{How do you test a platform, not a service?} ---
        Consumer-driven contract tests per surface \cite{pactdocs},
        in-process smoke suites proving surfaces coexist, conformance
        policy on every spec, and chaos drills on the integration points.
        (Ch.~\ref{ch:testing})
\end{enumerate}

\subsection*{Staff}

\begin{enumerate}[label=\textbf{Q\arabic*.}, leftmargin=3em, start=25]
  \item \textbf{Design the API strategy for a marketplace startup that will
        become a platform.} --- Year one: REST + OpenAPI public core,
        RFC~9457 and pagination conventions as day-one lint rules. Year two:
        add the profiles constraints demand (webhooks for partners, GraphQL
        if app payload graphs justify it), never by mandate. Year three:
        governance as code, rate tiers as a product, SLOs and error budgets,
        and a conformance checker gating every spec
        \cite{gehl2021apidesign}. (This chapter)
  \item \textbf{Walk me through your reference architecture's request
        lifecycle.} --- Edge (TLS, WAF, cache) $\rightarrow$ gateway (JWT
        verify, rate tier, correlation ID) $\rightarrow$ service (validate,
        idempotency check, gRPC inventory call, transaction + outbox)
        $\rightarrow$ relay $\rightarrow$ broker $\rightarrow$ signed webhook
        + SSE fan-out, with one OTel trace across all of it
        \cite{opentelemetrydocs}. (Section~\ref{sec:ref-arch})
  \item \textbf{How do MCP tool schemas change API design?} --- The consumer
        is now a model choosing tools by reading your schema: descriptions,
        naming, error structure, and idempotency semantics become behavioral
        surface area. Generate tool schemas from the OpenAPI source of truth
        \cite{openapi31,jsonschemadocs}, keep error shapes structured
        \cite{rfc9457}, and gate agent-callable surfaces behind the same
        authZ and audit discipline as human ones. (Section~\ref{sec:frontiers})
  \item \textbf{Your CEO read about HTTP/3 and wants it ``everywhere by
        Friday.'' Respond.} --- Public surfaces already inherit it at the CDN
        via \texttt{Alt-Svc} \cite{rfc9114}; east--west gRPC gains little and
        risks UDP filtering and CPU regressions; propose measuring p99 and
        CPU per request before and after an edge-only rollout.
        (Section~\ref{sec:frontiers})
  \item \textbf{How do you keep 50 teams' APIs coherent without a review
        board?} --- Executable policy: a handful of linted invariants
        (errors, auth, pagination, versioning) in CI, a self-service portal
        with templates that make the right way the easy way, and office
        hours for the genuinely hard cases \cite{openapi31,asyncapi}.
        (Ch.~\ref{ch:lifecycle})
  \item \textbf{What is your post-quantum readiness plan?} --- Inventory TLS
        termination points, track hybrid key-exchange support in your
        edge/library stack \cite{rfc8446}, automate certificate rotation so
        algorithm changes are routine, and watch handshake-size effects on
        QUIC amplification limits. (Section~\ref{sec:frontiers})
  \item \textbf{When is a second protocol stack worth its operational
        cost?} --- When a consumer class's constraints fail on the existing
        stack in a way money and caching cannot fix --- e.g., mobile
        over-fetching that no amount of BFF trimming solves, or a p99 that
        only binary framing meets. Quantify the constraint first; ``two
        stacks'' is a conclusion, never a starting point. (This chapter)
  \item \textbf{How do you evaluate an AI-generated API design?} --- The
        same way as a human one, plus distrust: run it through the decision
        card, lint it with executable policy, contract-test it against
        consumers, and treat its plausible-but-wrong inventions as the
        default case. AI drafts; deterministic gates decide.
        (Section~\ref{sec:frontiers})
\end{enumerate}

\section{External Bibliography \& Further Reading}

Consolidated across the book, organized by module. All entries are real,
primary, and freely available unless noted as books.

\subsection*{Module I --- Foundations}

\begin{itemize}
  \item Fielding's dissertation \cite{fielding2000rest} --- the REST
        constraints from the source; reread it after Chapter~\ref{ch:rest_style}.
  \item HTTP semantics and wire formats: RFC~9110 \cite{rfc9110},
        RFC~9112 \cite{rfc9112}, and for the modern transports RFC~9113
        \cite{rfc9113} with HPACK \cite{rfc7541}, and RFC~9114 over QUIC
        \cite{rfc9114,rfc9000}.
  \item JSON \cite{rfc8259} and JSON Schema \cite{jsonschemadocs} --- the
        serialization contract layer of Chapter~\ref{ch:serialization}.
  \item REST in practice: Richardson's maturity model book
        \cite{richardson2013restful}; API design craft in
        \cite{gehl2021apidesign}.
  \item Python HTTP clients: httpx \cite{httpxdocs} and requests
        \cite{requestsdocs}.
  \item TLS 1.3 \cite{rfc8446} for the transport security underpinning of
        Chapter~\ref{ch:client_security}.
\end{itemize}

\subsection*{Module II --- Building APIs}

\begin{itemize}
  \item FastAPI \cite{fastapidocs} and Pydantic \cite{pydanticdocs} ---
        the reference stack of Chapter~\ref{ch:fastapi}.
  \item RFC~9457 problem details \cite{rfc9457} --- the error interface of
        Chapter~\ref{ch:problem_details}.
  \item RFC~8288 Web Linking \cite{rfc8288} and RFC~8594 \texttt{Sunset}
        \cite{rfc8594} --- pagination links and deprecation signaling.
\end{itemize}

\subsection*{Module III --- Protocols \& Scale}

\begin{itemize}
  \item AsyncAPI \cite{asyncapi} --- the event-contract counterpart of
        OpenAPI for Chapter~\ref{ch:async_apis}.
  \item The GraphQL specification \cite{graphqlspec} --- read the execution
        section, not just the schema language.
  \item gRPC \cite{grpcdocs} and Protocol Buffers \cite{protobufdocs} ---
        including the versioning rules for \texttt{proto} files.
  \item The IETF RateLimit header draft \cite{ietf_ratelimit_headers} ---
        the emerging standard behind Chapter~\ref{ch:rate_limiting}.
  \item Kleppmann \cite{kleppmann2017ddia} (book) --- the data-systems
        foundation for outboxes, logs, and consistency reasoning.
  \item Nygard \cite{nygard2018releaseit} (book) --- the failure-mode
        catalog behind Chapter~\ref{ch:microservices}.
  \item Newman \cite{newman2021microservices} (book) --- service boundaries
        and decomposition strategy.
\end{itemize}

\subsection*{Module IV --- Production \& Governance}

\begin{itemize}
  \item OAuth~2.0 \cite{rfc6749}, PKCE \cite{rfc7636}, and JWT
        \cite{rfc7519} --- the token stack of
        Chapters~\ref{ch:client_security} and \ref{ch:api_security}.
  \item OWASP API Security Top 10 \cite{owaspapisecurity} --- the standing
        checklist; reread at every major release.
  \item OpenTelemetry \cite{opentelemetrydocs} and the Google SRE book
        \cite{beyer2016sre} --- instrumentation and SLO discipline for
        Chapter~\ref{ch:observability}.
  \item The twelve-factor methodology \cite{twelvefactor} --- configuration
        and deployment hygiene for Chapter~\ref{ch:deployment}.
  \item Pact \cite{pactdocs} --- consumer-driven contract testing for
        Chapter~\ref{ch:testing}.
  \item OpenAPI 3.1 \cite{openapi31} --- the governance source of truth for
        Chapter~\ref{ch:lifecycle}.
\end{itemize}

\section{Capstone Projects}
\label{sec:capstones}

Five integrative projects. Each builds on multiple chapters; datasets and
services are synthetic and offline.

\subsection{[junior] Project 1: Extend the Decision Advisor}

\textbf{Objective:} deepen Listing~1 (Section~\ref{sec:code-advisor}) into a usable team
tool.

\textbf{Inputs/Datasets:} \texttt{protocol\_advisor.py}, its test file, and
Table~\ref{tab:decision-criteria}.

\textbf{Steps:}
\begin{enumerate}
  \item Add an interactive questionnaire mode that prompts for the six
        criteria and prints the recommendation with its rationale trace.
  \item Add a seventh criterion (data-residency constraint) with signals
        that penalize public-CDN-dependent profiles when residency forbids
        edge caching.
  \item Add a \texttt{--explain} flag that prints the full score table for
        all five profiles, not only the winner.
  \item Extend the tests: the new criterion must change at least one
        canonical scenario's runner-up margin without flipping any existing
        winner.
\end{enumerate}

\textbf{Acceptance Criteria:}
\begin{itemize}
  \item[$\square$] Questionnaire mode runs offline and rejects invalid
        answers with a clear message.
  \item[$\square$] All pre-existing advisor tests pass unmodified.
  \item[$\square$] New tests cover the residency criterion and
        \texttt{--explain} output.
  \item[$\square$] Rationale traces mention the residency criterion when it
        moved a score.
\end{itemize}

\textbf{Stretch Goals:} emit the score table as JSON for CI consumption;
add a ``confidence'' metric based on the winner's margin.

\subsection{[mid] Project 2: Governance Toolchain}

\textbf{Objective:} grow the conformance checker of
Listing~3 (Section~\ref{sec:code-checker}) into a team-grade lint gate.

\textbf{Inputs/Datasets:} \texttt{conformance\_checker.py}, both fixture
folders, and the invariants of Chapter~\ref{ch:lifecycle}.

\textbf{Steps:}
\begin{enumerate}
  \item Add three new rules: every operation has an \texttt{operationId}
        matching \texttt{camelCase}; every error schema is registered in
        \texttt{components}, not inlined; AsyncAPI channel names follow
        \texttt{domain.noun.verb}.
  \item Produce both a console report and a machine-readable JSON report
        with stable rule codes.
  \item Add a \texttt{--fix-hints} mode that prints the smallest suggested
        remediation per violation.
  \item Wire the checker into a script that fails CI on any violation in a
        changed spec only (accept a folder pair: baseline vs.\ candidate).
\end{enumerate}

\textbf{Acceptance Criteria:}
\begin{itemize}
  \item[$\square$] All five original rule families plus the three new rules
        fire on \texttt{fixtures\_fail} and pass on \texttt{fixtures\_pass}.
  \item[$\square$] JSON report parses and round-trips deterministically.
  \item[$\square$] Changed-only mode ignores pre-existing violations in
        untouched files.
  \item[$\square$] Tests cover each new rule with a positive and negative
        fixture.
\end{itemize}

\textbf{Stretch Goals:} severity levels (error/warning) with per-rule
overrides in a policy file; SARIF output for code-scanning integration.

\subsection{[senior] Project 3: Multi-Protocol Read/Write Split}

\textbf{Objective:} build a domain where writes are REST, reads are GraphQL,
and internal coordination is gRPC --- over one shared store --- and prove
consistency.

\textbf{Inputs/Datasets:} Listing~2 (Section~\ref{sec:code-platform}) as the starting
skeleton; synthetic catalog/orders datasets from the Northwind universe.

\textbf{Steps:}
\begin{enumerate}
  \item Split Listing~2 (Section~\ref{sec:code-platform}) into a write service (REST,
        idempotent commands) and a read service (GraphQL projections).
  \item Propagate writes to the read store through an outbox-backed event
        stream (in-process broker fake is fine; keep the interface
        broker-shaped).
  \item Model the internal inventory check as a gRPC contract (write the
        \texttt{.proto}; an in-process stub implementation is acceptable
        offline).
  \item Write a consistency test: concurrent orders never oversell stock;
        the GraphQL read model converges within a bounded number of polls.
\end{enumerate}

\textbf{Acceptance Criteria:}
\begin{itemize}
  \item[$\square$] REST writes return RFC~9457 errors and honor idempotency
        keys.
  \item[$\square$] GraphQL reads never expose uncommitted state.
  \item[$\square$] The concurrency test passes deterministically in 100
        runs.
  \item[$\square$] OpenAPI and GraphQL SDL artifacts both exist and pass
        the Project-2 checker.
\end{itemize}

\textbf{Stretch Goals:} add a CDC-style lag metric to the read model and an
SLO on projection freshness; simulate broker outage and show recovery
without loss.

\subsection{[senior] Project 4: Webhook Notification Pipeline with Outbox Replay}

\textbf{Objective:} implement the partner-notification subsystem of the
reference architecture end to end.

\textbf{Inputs/Datasets:} the outbox pattern of
Chapter~\ref{ch:async_apis}; synthetic partner endpoint (a local FastAPI
receiver you also build).

\textbf{Steps:}
\begin{enumerate}
  \item Add an outbox table and relay loop to Listing~2 (Section~\ref{sec:code-platform})
        (in-process, polling, with exponential backoff).
  \item Sign every webhook with HMAC; build the receiving endpoint to verify
        signatures and reject replays outside a timestamp window.
  \item Persist a delivery log with attempt counts and last error; expose
        \texttt{POST /v1/webhooks/\{id\}/replay}.
  \item Emit AsyncAPI documentation for the event catalog; extend the
        conformance checker's A-001 to cover it.
  \item Add an SSE fallback stream carrying the same events for partners
        who cannot receive inbound webhooks.
\end{enumerate}

\textbf{Acceptance Criteria:}
\begin{itemize}
  \item[$\square$] No event is lost across a simulated broker/receiver
        outage (test proves it).
  \item[$\square$] Signature failures and stale timestamps are rejected
        with 401/422 problem details.
  \item[$\square$] Replay re-sends exactly the logged event payload.
  \item[$\square$] AsyncAPI spec passes the conformance checker.
\end{itemize}

\textbf{Stretch Goals:} dead-letter queue with operator triage endpoint;
delivery-latency histogram exported as an OTel-style metric.

\subsection{[staff] Project 5: The Grand Capstone --- Northwind Robotics API Platform}

\textbf{Objective:} build the complete platform of
Section~\ref{sec:grand-capstone}: all five surfaces of
Table~\ref{tab:capstone-surfaces}, the four milestones of
Table~\ref{tab:milestones}, and the acceptance rubric --- as a coherent,
governed, observable system.

\textbf{Inputs/Datasets:} the full chapter (reference architecture, NFRs,
milestone plan); synthetic Northwind datasets; in-process fakes for broker,
payments, and identity provider.

\textbf{Steps:}
\begin{enumerate}
  \item M1: REST catalog + orders core with keyset pagination, date
        versioning, RFC~9457 errors, idempotency keys, and contract tests.
  \item M2: GraphQL read model with depth limits and persisted queries;
        inventory gRPC service; gateway-simulating middleware for auth and
        rate tiers with IETF headers \cite{ietf_ratelimit_headers}.
  \item M3: outbox relay, signed webhooks with delivery log and replay
        (Project~4 folded in), SSE telemetry stream; OpenAPI + AsyncAPI
        manifests under conformance CI.
  \item M4: OTel-style trace context propagated across all surfaces; SLO
        definitions with a load test harness; chaos drill (kill inventory
        mid-order) demonstrating circuit-breaker behavior; developer-portal
        pages generated from the specs.
  \item Final: a written design review --- decision-card rationale per
        surface, frontier assessment for HTTP/3 and MCP, and a one-year
        evolution plan.
\end{enumerate}

\textbf{Acceptance Criteria:}
\begin{itemize}
  \item[$\square$] All five surfaces respond per spec in one offline
        in-process test run.
  \item[$\square$] Conformance checker passes on the full manifests folder.
  \item[$\square$] SLOs from Section~\ref{sec:grand-capstone} measured green
        against a recorded load test.
  \item[$\square$] Chaos drill produces the designed failure behavior with
        zero data loss (outbox replay verified).
  \item[$\square$] Rubric score $\geq$ 16/20 with no dimension below 3.
\end{itemize}

\textbf{Stretch Goals:} canary deployment simulation with traffic splitting;
MCP tool surface generated from the OpenAPI spec for the read-only catalog;
regional-affinity design document for a hypothetical EU expansion.

\section{Python Dependencies Appendix}

This chapter's code targets Python 3.11+ and pins minimum versions only,
per the book's convention:

\begin{minted}{text}
# requirements.txt
fastapi>=0.110
strawberry-graphql[fastapi]>=0.230
pydantic>=2.7
pyyaml>=6.0
httpx>=0.27
pytest>=8
\end{minted}

All installation commands below assume an activated virtual environment in
PowerShell (\texttt{.\textbackslash{}.venv\textbackslash{}Scripts\textbackslash{}Activate.ps1}).

\subsection{fastapi}

\textbf{Description:} the ASGI web framework used for every HTTP surface in
the platform slice \cite{fastapidocs}.

\textbf{Why included:} it is the book's reference framework
(Chapter~\ref{ch:fastapi}): type-driven validation, automatic OpenAPI
generation, and first-class streaming responses for SSE.

\textbf{Alternatives:} Flask (simpler, weaker typing), Django REST Framework
(heavier, ORM-coupled), Litestar (comparable, smaller ecosystem).

\textbf{Installation:}
\begin{minted}{powershell}
pip install "fastapi>=0.110"
\end{minted}

\subsection{strawberry-graphql[fastapi]}

\textbf{Description:} a code-first GraphQL library; the
\texttt{[fastapi]} extra provides the ASGI router used to mount the read
model at \texttt{/graphql}.

\textbf{Why included:} implements the GraphQL surface of the platform slice
with typed Python resolvers instead of a separate schema language
(Chapter~\ref{ch:graphql}); the extra keeps integration to one line.

\textbf{Alternatives:} Graphene (older API style), Ariadne (schema-first),
or hand-rolled JSON-over-POST (not recommended).

\textbf{Installation:}
\begin{minted}{powershell}
pip install "strawberry-graphql[fastapi]>=0.230"
\end{minted}

\subsection{pydantic}

\textbf{Description:} the data-validation library underlying FastAPI's
request/response models \cite{pydanticdocs}.

\textbf{Why included:} all request and response schemas in
Listing~2 (Section~\ref{sec:code-platform}) are Pydantic models; version 2's validator
performance and stricter typing matter at platform scale.

\textbf{Alternatives:} dataclasses + manual validation (verbose), msgspec
(faster, smaller ecosystem), attrs + cattrs (mature, more boilerplate).

\textbf{Installation:}
\begin{minted}{powershell}
pip install "pydantic>=2.7"
\end{minted}

\subsection{pyyaml}

\textbf{Description:} the YAML parser used by the conformance checker to
load OpenAPI and AsyncAPI manifests.

\textbf{Why included:} real-world API manifests are overwhelmingly YAML;
the checker needs nothing more than \texttt{yaml.safe\_load} plus structural
traversal, so a heavy spec parser would be over-engineering.

\textbf{Alternatives:} ruamel.yaml (round-trip preservation, heavier),
strictyaml (schema-restricted subset), or JSON-only manifests (loses
comments).

\textbf{Installation:}
\begin{minted}{powershell}
pip install "pyyaml>=6.0"
\end{minted}

\subsection{httpx}

\textbf{Description:} the modern sync/async HTTP client
\cite{httpxdocs}; pulled in by FastAPI's \texttt{TestClient} and used
throughout the book's client chapters.

\textbf{Why included:} the smoke suite's in-process client is built on it;
its transport abstraction is what makes network-free tests possible
(Chapter~\ref{ch:consuming_python}).

\textbf{Alternatives:} requests \cite{requestsdocs} (sync-only, no ASGI
transport), aiohttp (async-only, larger API surface).

\textbf{Installation:}
\begin{minted}{powershell}
pip install "httpx>=0.27"
\end{minted}

\subsection{pytest}

\textbf{Description:} the test runner and fixture framework for all three
listings' suites.

\textbf{Why included:} fixtures give each test an isolated platform
instance; \texttt{pytest.approx} and parametrization keep the suites terse;
it is the ecosystem default (Chapter~\ref{ch:testing}).

\textbf{Alternatives:} unittest (stdlib, more boilerplate), nose2 (fading),
doctest (documentation tests only).

\textbf{Installation:}
\begin{minted}{powershell}
pip install "pytest>=8"
\end{minted}

\section{Chapter Summary \& Key Takeaways}

\begin{itemize}
  \item Protocol selection is constraint satisfaction. Six criteria ---
        consumer, payload, latency, delivery, caching, team maturity ---
        determine the interaction profile; the Master Decision Card and its
        executable advisor make the choice auditable and teachable.
  \item Production platforms are multi-protocol by design. The Northwind
        reference architecture runs REST, GraphQL, gRPC, webhooks, and SSE
        simultaneously, unified by governed invariants rather than mandates.
  \item The transactional outbox, idempotency keys, and RFC~9457 errors are
        the three cheapest reliability investments in the entire book ---
        and the three most often deferred \cite{kleppmann2017ddia,rfc9457}.
  \item Governance is cheapest on day one and must be executable; a
        conformance checker turns architectural intent into CI reality
        \cite{openapi31,asyncapi}.
  \item Frontiers (HTTP/3, MCP, edge, post-quantum TLS, security automation)
        amplify existing maturity; adopt/trial/wait with explicit
        prerequisites, never enthusiasm alone.
  \item Mastery is a pathway: consume and build (junior), correctness
        (mid), scale and resilience (senior), strategy and governance
        (staff). The book's twenty-one chapters map one-to-one onto it.
\end{itemize}

\section{Quick-Reference Card}

\begin{table}[htbp]
  \centering
  \small
  \begin{tabularx}{\textwidth}{l X X}
    \toprule
    \multicolumn{3}{l}{\textbf{THE MASTER DECISION CARD --- pin this page.}} \\
    \midrule
    \textbf{If the constraint says\ldots} & \textbf{Then choose\ldots} & \textbf{And remember\ldots} \\
    \midrule
    unknown public consumers &
      REST + OpenAPI \cite{openapi31} &
      curl-ability is a feature; govern with lint from day one \\
    cacheable reads at scale &
      REST GETs \cite{rfc9110} &
      POST-based GraphQL bypasses every shared cache \\
    nested graphs, mobile/SPA &
      GraphQL \cite{graphqlspec} &
      depth limits + persisted queries are not optional \\
    internal, p99 $<$ 50\,ms &
      gRPC \cite{grpcdocs} &
      contract-first protos; version fields, never renumber \\
    fire-and-forget, machines &
      events / webhooks \cite{asyncapi} &
      transactional outbox or lose events silently \\
    server must push to browsers &
      SSE (WebSockets if bidirectional) &
      unary APIs cannot push --- structural, not preferential \\
    exploratory team &
      REST &
      maturity is a constraint, not an insult \\
    \midrule
    \textbf{Always} &
      RFC~9457 errors \cite{rfc9457}, idempotent writes, IETF rate headers
      \cite{ietf_ratelimit_headers}, OTel traces \cite{opentelemetrydocs},
      linted specs &
      these are protocol-independent invariants --- the actual platform \\
    \bottomrule
  \end{tabularx}
  \caption{The one-page Master Decision Card. Photocopy freely.}
  \label{tab:master-card}
\end{table}

\section{Closing Thoughts}

Twenty chapters ago, an API was a URL you fetched. It is now --- if this
book did its job --- a contract you design, a system you operate, a product
you govern, and a career you build deliberately. The tools will keep
changing: HTTP/3 will finish arriving, agents will become first-class
consumers, something not yet named will be next year's conference darling.
The disciplines will not change: say what you mean in a machine-readable
contract, fail in ways callers can handle, make retries safe, watch your
systems with honest telemetry, and govern invariants rather than tastes.
Engineers who internalize those sentences will never be behind the
frontier --- they will be the ones defining it. Go build platforms worth
depending on.

\section{Standardized Evidence Addendum}
\subsection{Benchmark Snapshot Template}
\begin{table}[h]
  \centering
  \small
  \begin{tabularx}{\textwidth}{l c c c c}
    \toprule
    \textbf{Config} & \textbf{p50 latency} & \textbf{p99 latency} & \textbf{Error rate} & \textbf{Throughput} \\
    \midrule
    Baseline & -- & -- & -- & 1.00x \\
    Candidate A & -- & -- & -- & -- \\
    Candidate B & -- & -- & -- & -- \\
    \bottomrule
  \end{tabularx}
  \caption{Minimum benchmark table to keep per chapter.}
\end{table}

\subsection{Request Trace Template}
\begin{itemize}
  \item \textbf{Request:} one representative API call for this chapter's topic.
  \item \textbf{Before:} behavior (headers, payload, latency, failure mode) with the naive approach.
  \item \textbf{After:} behavior with the technique in this chapter.
  \item \textbf{Result:} what changed in correctness, resilience, latency, and operability.
\end{itemize}

\subsection{Cost and Latency Budget Snapshot}
\begin{table}[h]
  \centering
  \small
  \begin{tabularx}{\textwidth}{l c c}
    \toprule
    \textbf{Stage} & \textbf{Latency budget} & \textbf{Cost driver} \\
    \midrule
    TLS / connection setup & -- & handshake round trips, keep-alive hit rate \\
    Request validation & -- & schema complexity, CPU per request \\
    Business logic and I/O & -- & downstream fan-out, query cost \\
    Serialization and egress & -- & payload size, compression ratio \\
    \bottomrule
  \end{tabularx}
  \caption{Operational budget template for this chapter's technique.}
\end{table}

\section{Configuration Preset Template}
\begin{minted}{text}
# api_policy.yaml
policy_version: "v1.0.0"
component: "<chapter_topic>"
timeout_ms: 0
retry_budget: 0
fallback_strategy: "fail_fast"
rollback_trigger: "error_budget_burn_gt_threshold"
\end{minted}

\section{CI Quality Gate Specification}
\begin{itemize}
  \item contract tests green against the pinned OpenAPI/AsyncAPI/Protobuf spec,
  \item no breaking schema changes without a version bump,
  \item error responses conform to RFC 9457 problem-details shape,
  \item benchmark deltas within approved regression bounds (latency and throughput),
  \item policy version and rollback plan present in release metadata.
\end{itemize}

\section{Incident Runbook Template}
\begin{enumerate}
  \item Detect regression from SLO burn-rate alerts or consumer reports.
  \item Scope impacted endpoints, versions, and consumer segments.
  \item Contain impact (rollback, feature flag, rate-limit tightening, or traffic shift).
  \item Diagnose with traces, structured logs, and contract diffs.
  \item Recover with validated fix and verified rollout procedure.
  \item Prevent recurrence via new regression tests and CI gates.
\end{enumerate}

\section{Cross-Chapter Decision Card}
\begin{table}[h]
  \centering
  \small
  \begin{tabularx}{\textwidth}{l X X}
    \toprule
    \textbf{Dimension} & \textbf{Choose this approach when} & \textbf{Prefer an alternative when} \\
    \midrule
    Use case & -- & -- \\
    Latency budget & -- & -- \\
    Operational cost & -- & -- \\
    Team maturity & -- & -- \\
    \bottomrule
  \end{tabularx}
  \caption{Decision card to complete per chapter.}
\end{table}

\section{ROI Model and Build-vs-Buy}
\begin{itemize}
  \item \textbf{Build when:} the capability is differentiating, compliance forbids vendors, or scale makes vendor pricing irrational.
  \item \textbf{Buy when:} the capability is commodity (auth, gateways, observability), time-to-market dominates, or staffing is thin.
  \item \textbf{Hidden costs to model:} on-call load, upgrade cadence, expertise ramp, data egress fees.
\end{itemize}

\section{Disaster Recovery Notes (RPO/RTO)}
\begin{itemize}
  \item Define recovery point objective (RPO): maximum acceptable data loss window.
  \item Define recovery time objective (RTO): maximum acceptable restoration time.
  \item Test restores regularly; an untested backup is not a backup.
\end{itemize}

